% appendix_playbook.tex — Vitrine Operator Playbook
% Practical user-guide appendix for the v5 report.
% Wire into main with:
%   \appendix
%   % appendix_playbook.tex — Vitrine Operator Playbook
% Practical user-guide appendix for the v5 report.
% Wire into main with:
%   \appendix
%   % appendix_playbook.tex — Vitrine Operator Playbook
% Practical user-guide appendix for the v5 report.
% Wire into main with:
%   \appendix
%   % appendix_playbook.tex — Vitrine Operator Playbook
% Practical user-guide appendix for the v5 report.
% Wire into main with:
%   \appendix
%   \input{v5/appendix_playbook}
%
% Dr John O'Hare, DreamLab AI Consulting Ltd — June 2026

% =====================================================================
\chapter{Operator Playbook: Running Vitrine End-to-End}
\label{app:playbook}
% =====================================================================

This appendix is a practical guide for a new operator.  It assumes the Docker
stack is accessible but does not assume familiarity with the internals; the
steps call the pipeline through its published Python and REST interfaces.
Read the capture section (\S\ref{sec:pb-capture}) before anything else:
\emph{capture quality is the single largest determinant of output quality},
and no downstream step can compensate for footage that was never sharp.

% ─────────────────────────────────────────────────────────────────────
\section{Capture Methodology: The Dominant Quality Lever}
\label{sec:pb-capture}
% ─────────────────────────────────────────────────────────────────────

\begin{keyfinding}
Motion blur is the dominant failure mode in practice.  Empirical measurement
on the \emph{dreamlab} dataset found zero under-observed area (median
$\approx$\,90 cameras per surface point) yet a holey, noisy reconstruction.
The problem was dense-but-motion-blurred footage, not insufficient coverage.
Blur cannot be recovered in post: selection can raise the sharpness floor, but
it cannot create detail the sensor never recorded.
\end{keyfinding}

\subsection{Camera settings}
\label{sec:pb-camera-settings}

Lock all automatic adjustments before pressing record.  Any automatic change
between frames is interpreted by COLMAP as a change in scene appearance and
degrades feature matching.

\begin{center}
\begin{tabular}{@{}L{3.2cm}L{4.8cm}L{5.8cm}@{}}
\toprule
\textbf{Setting} & \textbf{Recommendation} & \textbf{Why} \\
\midrule
Resolution     & 4K (3840$\times$2160)           & More features per frame; 8K adds storage cost without proportional gain as training downsamples anyway \\
Frame rate     & \textbf{60\,fps}                 & Caps exposure to $\leq$\,1/60\,s vs.\ 1/30\,s at 30\,fps, halving motion blur; also doubles the candidate frames for sharpness selection \\
Focus          & Manual, mid-distance (3--5\,m)  & Autofocus hunting softens frames and breaks feature matching \\
Exposure       & Manual, locked                  & Auto-exposure causes frame-to-frame brightness ramps that COLMAP interprets as surface-appearance changes \\
White balance  & Manual preset                   & Auto white balance shifts colour temperature between frames \\
Shutter speed  & $\geq$\,1/500\,s (lock in pro app) & Short exposure is the primary anti-blur lever; bright, even light lets you achieve this without under-exposure \\
Lens           & 24--50\,mm equivalent (1$\times$ main camera on phones) & Minimal distortion; ultra-wide ($<$16\,mm) fails in periphery; telephoto reduces parallax \\
\bottomrule
\end{tabular}
\end{center}

\medskip
\noindent\textbf{Phone-specific note.}  Flagship phones (iPhone 15 Pro,
Pixel 9 Pro, Samsung S24 Ultra) produce usable video \emph{only} if you use
the main 1$\times$ camera, lock exposure, focus, and white balance via a pro
app (e.g.\ mcpro24fps, FiLMiC Pro), and shoot in good light.  The Pixel 9 Pro
experiment found \textbf{4K/60\,fps} to be the best video mode: 60\,fps caps
exposure and doubles candidate frames.  \textbf{8K/30\,fps is worse}: the
longer exposure from 30\,fps produces blurrier frames; extra pixels on a
blurred frame are wasted.  If the camera supports RAW photo bursts, those are
even better --- very short exposure, no inter-frame compression.

\subsection{Movement and stabilisation}
\label{sec:pb-movement}

\begin{itemize}
  \item \textbf{Maximum walking speed: 0.5\,m/s.}  Counting ``one-Mississippi''
        per step is a practical governor.  When in doubt, slow down further.
  \item \textbf{Use a gimbal} (DJI RS, Zhiyun) or a tripod dolly.  These
        eliminate micro-shake and rolling-shutter artefacts.  Two-handed
        handheld capture is marginal in bright light; single-hand phone capture
        introduces motion blur on most frames and should be avoided.
  \item \textbf{Insert micro-pauses} at key viewpoints.  The sharpness
        sampler harvests the sharp still frames from these pauses.  Even a
        fraction of a second of steadiness at a corner or interesting object is
        worth more than continuous motion.
  \item \textbf{Move laterally, not just forward.}  The pipeline needs
        triangulation, and triangulation needs physical baseline between
        viewpoints.  Sidestep, arc, and weave rather than walking straight ahead.
  \item \textbf{Every point in the scene should be visible from at least five
        viewpoints.}  Revisit areas from different angles rather than passing
        through once.
\end{itemize}

\subsection{Capture patterns}
\label{sec:pb-patterns}

\paragraph{Room scan (60--120\,s per room).}
Walk the full perimeter at three height bands — low ($\approx$\,1\,m, waist),
eye level ($\approx$\,1.6\,m), and high (arms raised, $\approx$\,2.2\,m) —
then walk the two diagonals to fill the interior.  This gives three elevation
bands of every wall plus cross-coverage of floor and ceiling.  Multiple
independent sweeps of the same room at different heights combine into a denser
reconstruction than any single sweep: an observed three-sweep capture produced
$\approx$146\,k COLMAP points at $\approx$95\,\% registration versus $\approx$90\,k
at $\approx$90\,\% for a single perimeter pass.

\paragraph{Object orbit (30--60\,s per object).}
The best room sweep is rarely the best for any individual object.  Objects you
want as individual hulls need a \emph{dedicated} close orbit in addition to
the room sweeps:
\begin{enumerate}[nosep]
  \item Orbit at eye level ($\approx$1.5\,m distance): full 360°, 20\,s.
  \item Orbit at knee height, camera angled upward: full 360°, 20\,s.
  \item Orbit above head, camera angled downward: full 360°, 20\,s.
\end{enumerate}

\paragraph{Detail pass (10--15\,s per feature).}
After room and object passes, sweep 30--50\,cm from important surfaces (carved
detail, lettering, fine texture).  These close-up frames lift the resolution
ceiling for that region.

\subsection{Scene preparation}
\label{sec:pb-scene-prep}

\begin{itemize}
  \item \textbf{Remove all moving elements.}  People, animals, ceiling fans,
        screens playing video, swaying plants — any inter-frame motion creates
        smeared ghost geometry.  Pause recording if someone walks through the
        scene.  Even a person standing still shifts weight and breathes.
  \item \textbf{Turn screens off.}  A powered-off TV is far less glossy than a
        running one.  A running TV is the worst-case object: a large,
        view-dependent reflector whose reflections defeat both SfM (inconsistent
        features) and 3DGS (spherical-harmonics overfit the moving reflection
        into floaters).  The 2026-06-21 Drive dataset was abandoned for exactly
        this reason combined with low MUSIQ scores.
  \item \textbf{Cover mirrors and polished surfaces.}  Matte fabric or dulling
        spray prevents ghost geometry behind reflective surfaces.
  \item \textbf{Ensure even, diffuse lighting.}  Overcast sky is ideal outdoors.
        Indoors, use large soft panels or bounce lights off the ceiling.  Before
        recording, rotate 360° in the scene centre watching the exposure meter:
        if it swings more than 1.5 stops, add fill light or flag bright sources.
  \item \textbf{Add temporary texture to featureless surfaces.}  White walls and
        plain ceilings give COLMAP nothing to match.  Place books, plants, or
        patterned cloth against blank walls during capture; remove them later if
        needed.
\end{itemize}

\subsection{Do / do not summary}
\label{sec:pb-donot}

\begin{center}
\begin{tabular}{@{}L{5.8cm}L{7.2cm}@{}}
\toprule
\textbf{Do} & \textbf{Do not} \\
\midrule
4K/60\,fps, manual exposure, focus, WB & Use auto-exposure or auto white balance \\
Gimbal or two-handed stabilisation & Pan fast ($>$15°/s) or shoot single-handed \\
$\leq$\,0.5\,m/s, micro-pauses at key points & Walk at normal pace \\
Three height bands + diagonals & Walk only one straight path through the space \\
Dedicated 360° object orbits & Assume the room sweep covers individual objects \\
Transfer the original camera file over USB & Upload to a cloud service that transcodes \\
Even, diffuse lighting & Mix tungsten and daylight; allow hard shadows \\
Turn screens and fans off & Capture with any motion in the scene \\
$\geq$1/500\,s shutter (lock in pro app) & Rely on in-camera optical stabilisation alone for blur \\
4K/60\,fps on phones & Use 8K/30\,fps (blurrier despite higher pixel count) \\
\bottomrule
\end{tabular}
\end{center}

\subsection{Why blur is unrecoverable}
\label{sec:pb-blur-unrecoverable}

Motion blur is a measurement artefact: a blurred pixel is the average of many
spatial samples acquired over the exposure window, and the individual samples
are gone.  Per-frame generative deblurring (NAFNet, Restormer) hallucinates
texture rather than recovering it, and the hallucinated texture is spatially
inconsistent across frames, which breaks multi-view SfM.  The only
reconstruction methods that legitimately model blur --- joint exposure-trajectory
optimisation approaches such as \textsc{3dgs-deblur} --- show gains of
approximately 1--2\,dB PSNR on \emph{moderate} blur but hit the same
SfM-quality ceiling as the source footage.  \textbf{Blur is unrecoverable in
post; the real lever is short-shutter capture.}

The Vitrine frame-QA gate applies a MUSIQ no-reference IQA score plus a full-resolution
variance-of-Laplacian sharpness score to every extracted frame.  A video that cannot
supply \code{min\_good\_frames} above the KEEP threshold is marked
\code{too\_low\_quality} and the operator is asked to recapture.  The dreamlab
dataset scored a median MUSIQ of $\approx$\,31 (usable but blur-limited); the
abandoned Drive dataset scored $\approx$\,19 (unusable).  The gate is a diagnostic,
not a repair: it tells you when recapture is necessary.

% ─────────────────────────────────────────────────────────────────────
\section{Running the Pipeline}
\label{sec:pb-run}
% ─────────────────────────────────────────────────────────────────────

\subsection{Bringing up the stack}
\label{sec:pb-stack}

All pipeline work runs inside the \code{gaussian-toolkit} container.  From the
repository root on the host:

\begin{lstlisting}[language=bash,caption={Bring up the consolidated stack}]
# Clone with the vendored LichtFeld tool included
git clone --recurse-submodules <repo-url>
cd Vitrine

# Start the pipeline containers (GPU required)
docker compose -f docker-compose.consolidated.yml up -d

# Verify the stack is healthy
curl http://gaussian-toolkit:7860/health     # web UI
curl http://gaussian-toolkit:45677/mcp       # LichtFeld MCP
curl http://vitrine-comfyui:8188/system_stats  # ComfyUI (TRELLIS.2 / FLUX.2)
\end{lstlisting}

\noindent The Unreal Engine 5.8 overlay is optional and started separately:

\begin{lstlisting}[language=bash,caption={Bring up the Unreal overlay (optional)}]
docker compose \
  -f docker-compose.consolidated.yml \
  -f unreal/docker-compose.unreal.yml \
  up -d unreal unreal-mcp-bridge
\end{lstlisting}

\noindent Run the SOTA preflight check to verify model weights and VRAM before
processing any footage:

\begin{lstlisting}[language=bash,caption={SOTA preflight}]
python -m pipeline.sota_registry check
\end{lstlisting}

\noindent If any weight or pin is stale the preflight reports it; resolve before
continuing.  Model weights live in \code{data/comfyui/models/} ($\approx$\,216\,GB,
bind-mounted read-only); do not re-download them.

\subsection{Ingest and frame QA}
\label{sec:pb-ingest}

Place source video files in \code{data/raw/} and submit the job via the web UI
(\texttt{http://localhost:7860}) or directly:

\begin{lstlisting}[language=bash,caption={Step 1 -- ingest and frame extraction}]
python3 -c "
from pipeline.stages import PipelineStages
p = PipelineStages('/data/output/JOB_ID')
result = p.ingest('/data/raw/scene.mp4', fps=4.0)
print(result)
"
\end{lstlisting}

\noindent Ingest extracts frames at 4\,fps to oversample, then the frame-QA
stage selects the sharpest subset.  The QA pipeline applies two complementary
metrics:

\begin{enumerate}
  \item \textbf{MUSIQ no-reference IQA} (\code{pyiqa}, GPU).  A
        perceptual score on the KonIQ-10k scale ($\approx$0--100; good
        photography scores 50--75; severely blurred or noisy footage scores
        below $\approx$25).  Videos unable to supply enough frames above the
        KEEP threshold are marked \code{too\_low\_quality} and skipped.
  \item \textbf{Full-resolution variance-of-Laplacian.}  Sharpness is measured
        at full 4K resolution, not at a downscaled proxy --- downscaling blurs
        away the motion-blur signal (960p blurdetect showed only a
        $\approx$1.6$\times$ range on dreamlab frames; full-4K Laplacian showed a
        $\approx$105$\times$ range).  The \emph{sharpest frame per overlapping
        FIFO window} is kept, which raises the sharpness floor by 20--38\,\% at
        the bottom quartile --- the frames that most poison COLMAP matches.
\end{enumerate}

\noindent Inspect the selected frame count before proceeding:

\begin{lstlisting}[language=bash,caption={Step 3 -- select best frames}]
python3 -c "
from pipeline.stages import PipelineStages
p = PipelineStages('/data/output/JOB_ID')
result = p.select_frames('/data/output/JOB_ID/frames/', target=80)
print(result)
"
# Target: 60-80 frames. If < 40 are above threshold, consider recapture.
\end{lstlisting}

\subsection{Structure-from-Motion and 3DGS training}
\label{sec:pb-sfm-train}

\begin{lstlisting}[language=bash,caption={Step 4 -- COLMAP SfM (sequential matcher for video)}]
python3 -c "
from pipeline.stages import PipelineStages
p = PipelineStages('/data/output/JOB_ID')
result = p.reconstruct('/data/output/JOB_ID/frames_selected/',
                       matcher='sequential')
print(result)
"
# Quality gate: >= 70 % of frames must register.
# If < 50 % register, re-run with matcher='exhaustive'.
\end{lstlisting}

\begin{lstlisting}[language=bash,caption={Step 5 -- 3DGS training (30\,000 iterations)}]
python3 -c "
from pipeline.stages import PipelineStages
p = PipelineStages('/data/output/JOB_ID')
result = p.train('/data/output/JOB_ID/colmap/undistorted/',
                 iterations=30000)
print(result)
"
# Quality gate: PLY file should be > 10 MB; loss < 0.02.
\end{lstlisting}

\noindent The output splat is \code{output/JOB\_ID/model/splat\_30000.ply}.  This
file is the primary LichtFeld-native deliverable that feeds both the Gaussian
splat room path and the object segmentation branch.

\subsection{The stage sequence}
\label{sec:pb-sequence}

The full pipeline in order:

\begin{center}
\begin{tabular}{@{}lL{4.5cm}L{6cm}@{}}
\toprule
\textbf{Stage} & \textbf{What runs} & \textbf{Quality check} \\
\midrule
1 & Ingest: frame extraction at 4\,fps & \code{frames/} directory populated \\
2 & Remove people (if present) & Inspect sample frames \\
3 & Frame-QA selection & 60--80 KEEP-grade frames \\
4 & COLMAP SfM (sequential matcher) & $\geq$\,70\,\% registration \\
5 & 3DGS training (30\,k iterations) & PLY $>$10\,MB, loss $<$0.02 \\
6 & SAM3 concept segmentation & $\geq$\,2 objects identified \\
7 & Per-object PLY extraction + orbit renders & $\geq$5\,k vertices per object \\
7c & TRELLIS.2 hull\_e2e (textured GLB) & Textured GLB produced per object \\
8 & Room mesh extraction (TSDF or 2DGS) & $\geq$\,30\,k vertices, clean topology \\
8b & Mesh cleanup (\code{mesh\_cleaner smooth=0}) & No stray bridging triangles \\
8c & Texture bake (xatlas UV + vertex colours) & Texture atlas per mesh \\
9 & FBX export (\code{blender\_obj\_to\_fbx}) & FBX files produced \\
9b & UE 5.8 import (Nanite game assets) & Assets visible in Content Browser \\
10 & Validate + job complete & Full scene: room + objects \\
\bottomrule
\end{tabular}
\end{center}

\noindent \textbf{Do not stop after training.}  The pipeline is only complete
when it reaches the textured UE scene; stopping at the splat leaves the
deliverable unfinished.

% ─────────────────────────────────────────────────────────────────────
\section{Per-Capture Recipe Selection}
\label{sec:pb-recipe}
% ─────────────────────────────────────────────────────────────────────

Vitrine is not a fixed pipeline.  Each capture is \emph{diagnosed} first; the
router selects the configuration that best fits the capture's bottleneck.  The
four diagnostic axes are:

\begin{description}
  \item[Sharpness] Measured by MUSIQ + full-resolution Laplacian.  Is the
    footage blur-limited?
  \item[Coverage] Measured by a per-surface-point viewpoint count.  Are there
    under-observed regions (holes) or is coverage uniformly dense?
  \item[Hardware] RGB camera only, depth sensor (iPhone LiDAR), or event camera?
  \item[Deliverable] Fidelity-first splat room, polygonal mesh room, or both?
\end{description}

\noindent Table~\ref{tab:pb-recipes} maps common capture profiles to their
recommended configurations.

\begin{table}[ht]
\centering
\caption{Capture profile to recipe selector.  Objects are always reconstructed
via TRELLIS.2 \code{hull\_e2e} regardless of room profile.}
\label{tab:pb-recipes}
\begin{tabular}{@{}L{3.4cm}L{2.4cm}L{2.4cm}L{2.2cm}L{3.0cm}@{}}
\toprule
\textbf{Capture profile} & \textbf{Room rep} & \textbf{Enhancer} & \textbf{Mesh backend} & \textbf{Note} \\
\midrule
\textbf{Blurry $\cdot$ dense $\cdot$ RGB}
  (e.g.\ dreamlab)
  & Splat hero + mesh asset (dual)
  & \emph{None} (ArtiFixer won't help)
  & 2DGS / PGSR
  & Recapture is the real ceiling \\
\addlinespace
\textbf{Sharp $\cdot$ holey $\cdot$ RGB}
  & Splat (primary)
  & \textbf{ArtiFixer} (fills floaters/holes)
  & 2DGS / PGSR if mesh required
  & ArtiFixer's sweet spot \\
\addlinespace
\textbf{Sharp $\cdot$ dense $\cdot$ RGB}
  & Splat (primary)
  & None
  & 2DGS if mesh needed
  & Best case; least pain \\
\addlinespace
\textbf{iPhone LiDAR depth}
  & Mesh (primary)
  & Optional
  & \textbf{AGS-Mesh} (sensor depth)
  & Depth sensor unlocks clean walls \\
\addlinespace
\textbf{Event camera $\cdot$ blurry}
  & Splat
  & ---
  & ---
  & EBAD-Gaussian deblurs at source \\
\bottomrule
\end{tabular}
\end{table}

\subsection{Splat vs.\ mesh for the room}
\label{sec:pb-splat-vs-mesh}

\textbf{Default --- room as splat (NanoGS in UE 5.8).}  For an indoor room the
recommended default is to represent the room as a real Gaussian splat rendered
by the NanoGS UE 5.8 plugin, with object FBXs from TRELLIS.2 embedded
alongside it.  This is simultaneously the highest-fidelity path and the one
that leans most directly on the vendored LichtFeld trainer: the native
\code{splat\_30000.ply} is cleaned with SuperSplat, imported into UE 5.8, and
rendered as real Gaussians.  The mandated polygonal mesh is satisfied by the
TRELLIS.2 object FBXs.

\textbf{When to use a polygonal room mesh.}  If the brief requires a fully
polygonal game-asset scene (Nanite room mesh, not a splat), use \textbf{2DGS}
(quick win, same Open3D-TSDF backend as the current fallback, but 2D-surfel
depth is much cleaner than ellipsoid depth on flat walls) or \textbf{PGSR}
(planar Gaussians, explicitly handles textureless walls, strongest for indoor
rooms).  The legacy CoMe/Open3D-TSDF combination is retained as a fallback
but produces visibly blockier results and should not be the default.

\textbf{When ArtiFixer helps.}  ArtiFixer is a 3DGS refiner targeting
floaters, holes, and under-observed regions: it applies an opacity-mix
enhancement that fills in areas seen from too few viewpoints.  It runs on a
single 48\,GB RTX 6000 Ada (peak $\approx$\,45.5\,GB) and is the correct tool
for \emph{sharp but coverage-limited} captures.  It does \emph{not} deblur.
Applying ArtiFixer to a blur-limited capture wastes GPU time without improving
results.

\textbf{NanoGS floaters.}  The NanoGS plugin renders the splat faithfully,
including any floater Gaussians in the source PLY.  Pre-clean the PLY with
SuperSplat (floater removal by opacity and scale thresholds) before import.
If floaters persist after SuperSplat, check whether coverage was the issue
(ArtiFixer) or whether the source footage was simply blurry (recapture).

\subsection{Object reconstruction}
\label{sec:pb-objects}

Objects follow the same path regardless of room profile:

\begin{enumerate}
  \item \textbf{SAM3 concept segmentation} identifies freestanding objects in
        the splat and extracts per-object Gaussian PLY subsets.
  \item \textbf{Orbit render}: four synthetic camera views around the object
        centroid are rendered from the Gaussian subset.
  \item \textbf{TRELLIS.2 \code{hull\_e2e}}: the orbit renders are submitted to
        TRELLIS.2 (via ComfyUI at \code{http://vitrine-comfyui:8188}) which runs
        FLUX.2 turnaround generation followed by TRELLIS.2 3D reconstruction,
        producing a textured GLB with PBR materials.
  \item \textbf{Prescale}: Blender prescales the GLB to physical units to match
        the COLMAP coordinate frame.
  \item \textbf{FBX export}: the prescaled mesh is exported as FBX for UE 5.8
        game-asset import (Nanite-ready).
\end{enumerate}

\noindent TRELLIS.2 is the validated primary path; Hunyuan3D-2.1 (also staged
in ComfyUI) is the fallback for cases where the orbit renders are cluttered or
partially occluded.  Restart ComfyUI between objects if VRAM is under pressure:
FLUX.2 and TRELLIS.2 cannot co-reside with other large models on 48\,GB.

\subsection{VRAM phasing}
\label{sec:pb-vram}

The pipeline phases large models to avoid out-of-memory conditions.  The key
constraint is that FLUX.2 and TRELLIS.2 together require $\approx$\,45\,GB; they
cannot share VRAM with other resident models.  The orchestrator manages the
lifecycle automatically (load $\to$ infer $\to$ free) but the operator should
be aware:

\begin{itemize}
  \item If ComfyUI reports an OOM, the previous model was not freed cleanly.
        Restart the \code{vitrine-comfyui} container and re-run from the failing
        object.
  \item Process objects serially, not in parallel, unless a second GPU is
        dedicated to the task.
  \item ArtiFixer peaks at $\approx$\,45.5\,GB; do not run it concurrently with
        any ComfyUI workflow.
\end{itemize}

% ─────────────────────────────────────────────────────────────────────
\section{Troubleshooting}
\label{sec:pb-troubleshoot}
% ─────────────────────────────────────────────────────────────────────

\subsection{Recapture flag: footage too soft}
\label{sec:pb-ts-recapture}

\begin{technote}
\textbf{Symptom.}  The frame-QA stage returns \code{VideoQualityVerdict.too\_low\_quality} for all
or most input videos.  The MUSIQ median is below $\approx$\,25.
\end{technote}

The QA gate cannot manufacture sharpness that was never recorded.  Lowering
the MUSIQ threshold feeds more frames to COLMAP but does not improve the
reconstruction --- it just postpones a failure that COLMAP will surface as low
registration or that 3DGS training will surface as a blurry, noisy splat.

\textbf{Resolution: recapture} following the camera settings and movement
guidance in \S\ref{sec:pb-camera-settings}--\ref{sec:pb-movement}.  Checklist:
\begin{itemize}[nosep]
  \item Switch to 60\,fps (halves exposure time vs.\ 30\,fps).
  \item Lock shutter at $\geq$\,1/500\,s in a pro camera app.
  \item Reduce movement speed; add micro-pauses at key viewpoints.
  \item Use a gimbal or brace the camera firmly.
  \item Ensure bright, even illumination (adequate light allows fast shutter
        without under-exposure).
\end{itemize}

\subsection{Directional blur gate}
\label{sec:pb-ts-blur-gate}

Vitrine applies a directional blur metric in addition to the MUSIQ score:
optical-flow magnitude and direction identify frames where camera motion was
predominantly linear (walking forward) versus rotational (panning).  Panning
motion at low shutter speeds produces pronounced horizontal or vertical blur
streaks that fool the Laplacian sharpness metric (which measures spatial
frequency without direction).  Frames with high directional blur are downgraded
before FIFO windowed selection.

If the frame-QA gate discards a high fraction of frames despite the footage
appearing ``mostly sharp'', check whether long pan segments are responsible:
inspect the flow field overlays saved to \code{output/JOB\_ID/qa\_debug/}.
Recapture with slower, more deliberate panning ($<$\,15°/s) or replace panning
segments with lateral translation.

\subsection{COLMAP: low registration rate}
\label{sec:pb-ts-colmap}

A COLMAP registration below 70\,\% usually indicates one of the following:

\begin{description}
  \item[Insufficient overlap] You moved too fast or the scene was only captured
    from one direction.  Recapture with slower movement and the systematic
    multi-pass patterns described in \S\ref{sec:pb-patterns}.
  \item[COLMAP invocation flags] With COLMAP 4.1.0 the option namespace changed;
    old recipes silently use wrong flags.  Use \code{--FeatureExtraction.*} and
    \code{--FeatureMatching.*} (not the old \code{--SiftExtraction.*} namespace).
    Run feature extraction on GPU (\code{--FeatureExtraction.type SIFT} on GPU):
    GPU SIFT achieved $\approx$\,100\,\% registration on the multi-sweep
    dreamlab captures.
  \item[Missing output directory] COLMAP silently fails to create its database if
    the output directory does not already exist.  Run
    \code{mkdir -p colmap/} before invoking \code{feature\_extractor}.
  \item[Mixed cameras] If the capture combined multiple devices, pass
    \code{--ImageReader.single\_camera\_per\_folder} so COLMAP does not force
    one shared intrinsic model across devices.
  \item[Featureless surfaces] Large white walls or uniform floors provide no
    feature points.  Add temporary texture to blank surfaces before
    recapturing (\S\ref{sec:pb-scene-prep}).
\end{description}

\subsection{NanoGS floaters in the UE scene}
\label{sec:pb-ts-floaters}

\begin{description}
  \item[Cause: coverage holes.]  Regions seen from fewer than five viewpoints
    generate noisy or un-anchored Gaussians.  Apply ArtiFixer (the correct tool
    for this profile) before running SuperSplat.
  \item[Cause: motion blur in source footage.]  A blurry Gaussian is still a
    Gaussian; NanoGS renders it faithfully as a blurry blob.  ArtiFixer does
    not help here.  The resolution is recapture.
  \item[Prune with SuperSplat before UE import.] SuperSplat's opacity and scale
    thresholds reliably remove stray low-opacity Gaussians that accumulate in
    poorly observed regions.  Run SuperSplat on every PLY before importing into
    UE regardless of whether ArtiFixer was applied.
\end{description}

\subsection{UE import scale and object placement}
\label{sec:pb-ts-ue-scale}

UE 5.8 imports FBX in centimetres by default; the COLMAP/3DGS coordinate
frame is in metres.  Vitrine's \code{blender\_prescale\_objects} step converts
object meshes to physical scale before FBX export.  If imported objects appear
100$\times$ too large or too small, check that \code{blender\_prescale\_objects}
ran successfully and that the UE import dialogue has \textit{Import Uniform
Scale} set to 1.0 (not 0.01 or 100).  Use \code{ue\_place\_prescaled} (the
keeper placement script) rather than hand-placing objects via Remote Control,
which has been observed to produce incorrect scale on the final assembly step.

\subsection{GPU VRAM OOM during ComfyUI object processing}
\label{sec:pb-ts-vram}

FLUX.2 and TRELLIS.2 together peak at $\approx$\,45.5\,GB.  If a workflow
fails with an out-of-memory error:
\begin{enumerate}[nosep]
  \item Restart the \code{vitrine-comfyui} container to flush all resident
        models: \code{docker restart vitrine-comfyui}.
  \item Re-submit the failing object workflow.  ComfyUI's queue is persistent
        across restarts; the job will resume from the failed node.
  \item If OOM recurs, check that ArtiFixer is not concurrently resident.
        ArtiFixer should be run as a separate stage with ComfyUI idle.
  \item If the \code{vitrine-comfyui} container is sharing GPU\,0 with
        \code{gaussian-toolkit} 3DGS training, defer object processing until
        training is complete.
\end{enumerate}

\noindent The \code{--local\_attn\_size} flag in ArtiFixer is the dominant
VRAM lever: reducing it from the default to a smaller window size (e.g.\ 128)
significantly reduces peak usage at a modest quality cost.  See
\code{docker/artifixer/patches/} for the confirmed fix applied to the
dreamlab run.

% ─────────────────────────────────────────────────────────────────────
\section{Pre-Capture Checklist}
\label{sec:pb-checklist}
% ─────────────────────────────────────────────────────────────────────

Print and check before every capture session.

\begin{center}
\fbox{\begin{minipage}{0.92\textwidth}
\textbf{Pre-capture}\\[2pt]
$\square$\ Scene is fully static: no people, pets, fans, active screens, or flowing water\\
$\square$\ Lighting is even and consistent (meter swing $<$\,1.5 stops over 360°)\\
$\square$\ Mirrors and reflective surfaces covered or removed\\
$\square$\ Featureless walls have temporary texture added\\
$\square$\ Camera battery charged; storage space verified\\[6pt]
\textbf{Camera settings}\\[2pt]
$\square$\ Resolution: 4K; frame rate: 60\,fps\\
$\square$\ Focus: MANUAL, mid-distance\\
$\square$\ Exposure: MANUAL, locked for average scene brightness\\
$\square$\ Shutter: $\geq$\,1/500\,s (locked in pro app)\\
$\square$\ White balance: MANUAL preset matched to light source\\
$\square$\ Lens: 24--50\,mm equivalent (1$\times$ main camera on phones)\\[6pt]
\textbf{Capture}\\[2pt]
$\square$\ Using gimbal or braced two-handed stabilisation\\
$\square$\ Movement speed: $\leq$\,0.5\,m/s; micro-pauses at key viewpoints\\
$\square$\ Perimeter passes at three heights (low / eye / high)\\
$\square$\ Centre diagonal cross-passes\\
$\square$\ Dedicated 360° orbit of each key object at three heights\\
$\square$\ Detail pass of important features (30--50\,cm, 10--15\,s each)\\
$\square$\ Total duration: 60--120\,s per room; 30--60\,s per object\\[6pt]
\textbf{Post-capture}\\[2pt]
$\square$\ Transfer original file via USB or AirDrop (not cloud)\\
$\square$\ No re-encoding unless unavoidable; if required: H.264 CRF\,18, original resolution\\
$\square$\ File placed in \code{data/raw/} on the pipeline host\\
\end{minipage}}
\end{center}

%
% Dr John O'Hare, DreamLab AI Consulting Ltd — June 2026

% =====================================================================
\chapter{Operator Playbook: Running Vitrine End-to-End}
\label{app:playbook}
% =====================================================================

This appendix is a practical guide for a new operator.  It assumes the Docker
stack is accessible but does not assume familiarity with the internals; the
steps call the pipeline through its published Python and REST interfaces.
Read the capture section (\S\ref{sec:pb-capture}) before anything else:
\emph{capture quality is the single largest determinant of output quality},
and no downstream step can compensate for footage that was never sharp.

% ─────────────────────────────────────────────────────────────────────
\section{Capture Methodology: The Dominant Quality Lever}
\label{sec:pb-capture}
% ─────────────────────────────────────────────────────────────────────

\begin{keyfinding}
Motion blur is the dominant failure mode in practice.  Empirical measurement
on the \emph{dreamlab} dataset found zero under-observed area (median
$\approx$\,90 cameras per surface point) yet a holey, noisy reconstruction.
The problem was dense-but-motion-blurred footage, not insufficient coverage.
Blur cannot be recovered in post: selection can raise the sharpness floor, but
it cannot create detail the sensor never recorded.
\end{keyfinding}

\subsection{Camera settings}
\label{sec:pb-camera-settings}

Lock all automatic adjustments before pressing record.  Any automatic change
between frames is interpreted by COLMAP as a change in scene appearance and
degrades feature matching.

\begin{center}
\begin{tabular}{@{}L{3.2cm}L{4.8cm}L{5.8cm}@{}}
\toprule
\textbf{Setting} & \textbf{Recommendation} & \textbf{Why} \\
\midrule
Resolution     & 4K (3840$\times$2160)           & More features per frame; 8K adds storage cost without proportional gain as training downsamples anyway \\
Frame rate     & \textbf{60\,fps}                 & Caps exposure to $\leq$\,1/60\,s vs.\ 1/30\,s at 30\,fps, halving motion blur; also doubles the candidate frames for sharpness selection \\
Focus          & Manual, mid-distance (3--5\,m)  & Autofocus hunting softens frames and breaks feature matching \\
Exposure       & Manual, locked                  & Auto-exposure causes frame-to-frame brightness ramps that COLMAP interprets as surface-appearance changes \\
White balance  & Manual preset                   & Auto white balance shifts colour temperature between frames \\
Shutter speed  & $\geq$\,1/500\,s (lock in pro app) & Short exposure is the primary anti-blur lever; bright, even light lets you achieve this without under-exposure \\
Lens           & 24--50\,mm equivalent (1$\times$ main camera on phones) & Minimal distortion; ultra-wide ($<$16\,mm) fails in periphery; telephoto reduces parallax \\
\bottomrule
\end{tabular}
\end{center}

\medskip
\noindent\textbf{Phone-specific note.}  Flagship phones (iPhone 15 Pro,
Pixel 9 Pro, Samsung S24 Ultra) produce usable video \emph{only} if you use
the main 1$\times$ camera, lock exposure, focus, and white balance via a pro
app (e.g.\ mcpro24fps, FiLMiC Pro), and shoot in good light.  The Pixel 9 Pro
experiment found \textbf{4K/60\,fps} to be the best video mode: 60\,fps caps
exposure and doubles candidate frames.  \textbf{8K/30\,fps is worse}: the
longer exposure from 30\,fps produces blurrier frames; extra pixels on a
blurred frame are wasted.  If the camera supports RAW photo bursts, those are
even better --- very short exposure, no inter-frame compression.

\subsection{Movement and stabilisation}
\label{sec:pb-movement}

\begin{itemize}
  \item \textbf{Maximum walking speed: 0.5\,m/s.}  Counting ``one-Mississippi''
        per step is a practical governor.  When in doubt, slow down further.
  \item \textbf{Use a gimbal} (DJI RS, Zhiyun) or a tripod dolly.  These
        eliminate micro-shake and rolling-shutter artefacts.  Two-handed
        handheld capture is marginal in bright light; single-hand phone capture
        introduces motion blur on most frames and should be avoided.
  \item \textbf{Insert micro-pauses} at key viewpoints.  The sharpness
        sampler harvests the sharp still frames from these pauses.  Even a
        fraction of a second of steadiness at a corner or interesting object is
        worth more than continuous motion.
  \item \textbf{Move laterally, not just forward.}  The pipeline needs
        triangulation, and triangulation needs physical baseline between
        viewpoints.  Sidestep, arc, and weave rather than walking straight ahead.
  \item \textbf{Every point in the scene should be visible from at least five
        viewpoints.}  Revisit areas from different angles rather than passing
        through once.
\end{itemize}

\subsection{Capture patterns}
\label{sec:pb-patterns}

\paragraph{Room scan (60--120\,s per room).}
Walk the full perimeter at three height bands — low ($\approx$\,1\,m, waist),
eye level ($\approx$\,1.6\,m), and high (arms raised, $\approx$\,2.2\,m) —
then walk the two diagonals to fill the interior.  This gives three elevation
bands of every wall plus cross-coverage of floor and ceiling.  Multiple
independent sweeps of the same room at different heights combine into a denser
reconstruction than any single sweep: an observed three-sweep capture produced
$\approx$146\,k COLMAP points at $\approx$95\,\% registration versus $\approx$90\,k
at $\approx$90\,\% for a single perimeter pass.

\paragraph{Object orbit (30--60\,s per object).}
The best room sweep is rarely the best for any individual object.  Objects you
want as individual hulls need a \emph{dedicated} close orbit in addition to
the room sweeps:
\begin{enumerate}[nosep]
  \item Orbit at eye level ($\approx$1.5\,m distance): full 360°, 20\,s.
  \item Orbit at knee height, camera angled upward: full 360°, 20\,s.
  \item Orbit above head, camera angled downward: full 360°, 20\,s.
\end{enumerate}

\paragraph{Detail pass (10--15\,s per feature).}
After room and object passes, sweep 30--50\,cm from important surfaces (carved
detail, lettering, fine texture).  These close-up frames lift the resolution
ceiling for that region.

\subsection{Scene preparation}
\label{sec:pb-scene-prep}

\begin{itemize}
  \item \textbf{Remove all moving elements.}  People, animals, ceiling fans,
        screens playing video, swaying plants — any inter-frame motion creates
        smeared ghost geometry.  Pause recording if someone walks through the
        scene.  Even a person standing still shifts weight and breathes.
  \item \textbf{Turn screens off.}  A powered-off TV is far less glossy than a
        running one.  A running TV is the worst-case object: a large,
        view-dependent reflector whose reflections defeat both SfM (inconsistent
        features) and 3DGS (spherical-harmonics overfit the moving reflection
        into floaters).  The 2026-06-21 Drive dataset was abandoned for exactly
        this reason combined with low MUSIQ scores.
  \item \textbf{Cover mirrors and polished surfaces.}  Matte fabric or dulling
        spray prevents ghost geometry behind reflective surfaces.
  \item \textbf{Ensure even, diffuse lighting.}  Overcast sky is ideal outdoors.
        Indoors, use large soft panels or bounce lights off the ceiling.  Before
        recording, rotate 360° in the scene centre watching the exposure meter:
        if it swings more than 1.5 stops, add fill light or flag bright sources.
  \item \textbf{Add temporary texture to featureless surfaces.}  White walls and
        plain ceilings give COLMAP nothing to match.  Place books, plants, or
        patterned cloth against blank walls during capture; remove them later if
        needed.
\end{itemize}

\subsection{Do / do not summary}
\label{sec:pb-donot}

\begin{center}
\begin{tabular}{@{}L{5.8cm}L{7.2cm}@{}}
\toprule
\textbf{Do} & \textbf{Do not} \\
\midrule
4K/60\,fps, manual exposure, focus, WB & Use auto-exposure or auto white balance \\
Gimbal or two-handed stabilisation & Pan fast ($>$15°/s) or shoot single-handed \\
$\leq$\,0.5\,m/s, micro-pauses at key points & Walk at normal pace \\
Three height bands + diagonals & Walk only one straight path through the space \\
Dedicated 360° object orbits & Assume the room sweep covers individual objects \\
Transfer the original camera file over USB & Upload to a cloud service that transcodes \\
Even, diffuse lighting & Mix tungsten and daylight; allow hard shadows \\
Turn screens and fans off & Capture with any motion in the scene \\
$\geq$1/500\,s shutter (lock in pro app) & Rely on in-camera optical stabilisation alone for blur \\
4K/60\,fps on phones & Use 8K/30\,fps (blurrier despite higher pixel count) \\
\bottomrule
\end{tabular}
\end{center}

\subsection{Why blur is unrecoverable}
\label{sec:pb-blur-unrecoverable}

Motion blur is a measurement artefact: a blurred pixel is the average of many
spatial samples acquired over the exposure window, and the individual samples
are gone.  Per-frame generative deblurring (NAFNet, Restormer) hallucinates
texture rather than recovering it, and the hallucinated texture is spatially
inconsistent across frames, which breaks multi-view SfM.  The only
reconstruction methods that legitimately model blur --- joint exposure-trajectory
optimisation approaches such as \textsc{3dgs-deblur} --- show gains of
approximately 1--2\,dB PSNR on \emph{moderate} blur but hit the same
SfM-quality ceiling as the source footage.  \textbf{Blur is unrecoverable in
post; the real lever is short-shutter capture.}

The Vitrine frame-QA gate applies a MUSIQ no-reference IQA score plus a full-resolution
variance-of-Laplacian sharpness score to every extracted frame.  A video that cannot
supply \code{min\_good\_frames} above the KEEP threshold is marked
\code{too\_low\_quality} and the operator is asked to recapture.  The dreamlab
dataset scored a median MUSIQ of $\approx$\,31 (usable but blur-limited); the
abandoned Drive dataset scored $\approx$\,19 (unusable).  The gate is a diagnostic,
not a repair: it tells you when recapture is necessary.

% ─────────────────────────────────────────────────────────────────────
\section{Running the Pipeline}
\label{sec:pb-run}
% ─────────────────────────────────────────────────────────────────────

\subsection{Bringing up the stack}
\label{sec:pb-stack}

All pipeline work runs inside the \code{gaussian-toolkit} container.  From the
repository root on the host:

\begin{lstlisting}[language=bash,caption={Bring up the consolidated stack}]
# Clone with the vendored LichtFeld tool included
git clone --recurse-submodules <repo-url>
cd Vitrine

# Start the pipeline containers (GPU required)
docker compose -f docker-compose.consolidated.yml up -d

# Verify the stack is healthy
curl http://gaussian-toolkit:7860/health     # web UI
curl http://gaussian-toolkit:45677/mcp       # LichtFeld MCP
curl http://vitrine-comfyui:8188/system_stats  # ComfyUI (TRELLIS.2 / FLUX.2)
\end{lstlisting}

\noindent The Unreal Engine 5.8 overlay is optional and started separately:

\begin{lstlisting}[language=bash,caption={Bring up the Unreal overlay (optional)}]
docker compose \
  -f docker-compose.consolidated.yml \
  -f unreal/docker-compose.unreal.yml \
  up -d unreal unreal-mcp-bridge
\end{lstlisting}

\noindent Run the SOTA preflight check to verify model weights and VRAM before
processing any footage:

\begin{lstlisting}[language=bash,caption={SOTA preflight}]
python -m pipeline.sota_registry check
\end{lstlisting}

\noindent If any weight or pin is stale the preflight reports it; resolve before
continuing.  Model weights live in \code{data/comfyui/models/} ($\approx$\,216\,GB,
bind-mounted read-only); do not re-download them.

\subsection{Ingest and frame QA}
\label{sec:pb-ingest}

Place source video files in \code{data/raw/} and submit the job via the web UI
(\texttt{http://localhost:7860}) or directly:

\begin{lstlisting}[language=bash,caption={Step 1 -- ingest and frame extraction}]
python3 -c "
from pipeline.stages import PipelineStages
p = PipelineStages('/data/output/JOB_ID')
result = p.ingest('/data/raw/scene.mp4', fps=4.0)
print(result)
"
\end{lstlisting}

\noindent Ingest extracts frames at 4\,fps to oversample, then the frame-QA
stage selects the sharpest subset.  The QA pipeline applies two complementary
metrics:

\begin{enumerate}
  \item \textbf{MUSIQ no-reference IQA} (\code{pyiqa}, GPU).  A
        perceptual score on the KonIQ-10k scale ($\approx$0--100; good
        photography scores 50--75; severely blurred or noisy footage scores
        below $\approx$25).  Videos unable to supply enough frames above the
        KEEP threshold are marked \code{too\_low\_quality} and skipped.
  \item \textbf{Full-resolution variance-of-Laplacian.}  Sharpness is measured
        at full 4K resolution, not at a downscaled proxy --- downscaling blurs
        away the motion-blur signal (960p blurdetect showed only a
        $\approx$1.6$\times$ range on dreamlab frames; full-4K Laplacian showed a
        $\approx$105$\times$ range).  The \emph{sharpest frame per overlapping
        FIFO window} is kept, which raises the sharpness floor by 20--38\,\% at
        the bottom quartile --- the frames that most poison COLMAP matches.
\end{enumerate}

\noindent Inspect the selected frame count before proceeding:

\begin{lstlisting}[language=bash,caption={Step 3 -- select best frames}]
python3 -c "
from pipeline.stages import PipelineStages
p = PipelineStages('/data/output/JOB_ID')
result = p.select_frames('/data/output/JOB_ID/frames/', target=80)
print(result)
"
# Target: 60-80 frames. If < 40 are above threshold, consider recapture.
\end{lstlisting}

\subsection{Structure-from-Motion and 3DGS training}
\label{sec:pb-sfm-train}

\begin{lstlisting}[language=bash,caption={Step 4 -- COLMAP SfM (sequential matcher for video)}]
python3 -c "
from pipeline.stages import PipelineStages
p = PipelineStages('/data/output/JOB_ID')
result = p.reconstruct('/data/output/JOB_ID/frames_selected/',
                       matcher='sequential')
print(result)
"
# Quality gate: >= 70 % of frames must register.
# If < 50 % register, re-run with matcher='exhaustive'.
\end{lstlisting}

\begin{lstlisting}[language=bash,caption={Step 5 -- 3DGS training (30\,000 iterations)}]
python3 -c "
from pipeline.stages import PipelineStages
p = PipelineStages('/data/output/JOB_ID')
result = p.train('/data/output/JOB_ID/colmap/undistorted/',
                 iterations=30000)
print(result)
"
# Quality gate: PLY file should be > 10 MB; loss < 0.02.
\end{lstlisting}

\noindent The output splat is \code{output/JOB\_ID/model/splat\_30000.ply}.  This
file is the primary LichtFeld-native deliverable that feeds both the Gaussian
splat room path and the object segmentation branch.

\subsection{The stage sequence}
\label{sec:pb-sequence}

The full pipeline in order:

\begin{center}
\begin{tabular}{@{}lL{4.5cm}L{6cm}@{}}
\toprule
\textbf{Stage} & \textbf{What runs} & \textbf{Quality check} \\
\midrule
1 & Ingest: frame extraction at 4\,fps & \code{frames/} directory populated \\
2 & Remove people (if present) & Inspect sample frames \\
3 & Frame-QA selection & 60--80 KEEP-grade frames \\
4 & COLMAP SfM (sequential matcher) & $\geq$\,70\,\% registration \\
5 & 3DGS training (30\,k iterations) & PLY $>$10\,MB, loss $<$0.02 \\
6 & SAM3 concept segmentation & $\geq$\,2 objects identified \\
7 & Per-object PLY extraction + orbit renders & $\geq$5\,k vertices per object \\
7c & TRELLIS.2 hull\_e2e (textured GLB) & Textured GLB produced per object \\
8 & Room mesh extraction (TSDF or 2DGS) & $\geq$\,30\,k vertices, clean topology \\
8b & Mesh cleanup (\code{mesh\_cleaner smooth=0}) & No stray bridging triangles \\
8c & Texture bake (xatlas UV + vertex colours) & Texture atlas per mesh \\
9 & FBX export (\code{blender\_obj\_to\_fbx}) & FBX files produced \\
9b & UE 5.8 import (Nanite game assets) & Assets visible in Content Browser \\
10 & Validate + job complete & Full scene: room + objects \\
\bottomrule
\end{tabular}
\end{center}

\noindent \textbf{Do not stop after training.}  The pipeline is only complete
when it reaches the textured UE scene; stopping at the splat leaves the
deliverable unfinished.

% ─────────────────────────────────────────────────────────────────────
\section{Per-Capture Recipe Selection}
\label{sec:pb-recipe}
% ─────────────────────────────────────────────────────────────────────

Vitrine is not a fixed pipeline.  Each capture is \emph{diagnosed} first; the
router selects the configuration that best fits the capture's bottleneck.  The
four diagnostic axes are:

\begin{description}
  \item[Sharpness] Measured by MUSIQ + full-resolution Laplacian.  Is the
    footage blur-limited?
  \item[Coverage] Measured by a per-surface-point viewpoint count.  Are there
    under-observed regions (holes) or is coverage uniformly dense?
  \item[Hardware] RGB camera only, depth sensor (iPhone LiDAR), or event camera?
  \item[Deliverable] Fidelity-first splat room, polygonal mesh room, or both?
\end{description}

\noindent Table~\ref{tab:pb-recipes} maps common capture profiles to their
recommended configurations.

\begin{table}[ht]
\centering
\caption{Capture profile to recipe selector.  Objects are always reconstructed
via TRELLIS.2 \code{hull\_e2e} regardless of room profile.}
\label{tab:pb-recipes}
\begin{tabular}{@{}L{3.4cm}L{2.4cm}L{2.4cm}L{2.2cm}L{3.0cm}@{}}
\toprule
\textbf{Capture profile} & \textbf{Room rep} & \textbf{Enhancer} & \textbf{Mesh backend} & \textbf{Note} \\
\midrule
\textbf{Blurry $\cdot$ dense $\cdot$ RGB}
  (e.g.\ dreamlab)
  & Splat hero + mesh asset (dual)
  & \emph{None} (ArtiFixer won't help)
  & 2DGS / PGSR
  & Recapture is the real ceiling \\
\addlinespace
\textbf{Sharp $\cdot$ holey $\cdot$ RGB}
  & Splat (primary)
  & \textbf{ArtiFixer} (fills floaters/holes)
  & 2DGS / PGSR if mesh required
  & ArtiFixer's sweet spot \\
\addlinespace
\textbf{Sharp $\cdot$ dense $\cdot$ RGB}
  & Splat (primary)
  & None
  & 2DGS if mesh needed
  & Best case; least pain \\
\addlinespace
\textbf{iPhone LiDAR depth}
  & Mesh (primary)
  & Optional
  & \textbf{AGS-Mesh} (sensor depth)
  & Depth sensor unlocks clean walls \\
\addlinespace
\textbf{Event camera $\cdot$ blurry}
  & Splat
  & ---
  & ---
  & EBAD-Gaussian deblurs at source \\
\bottomrule
\end{tabular}
\end{table}

\subsection{Splat vs.\ mesh for the room}
\label{sec:pb-splat-vs-mesh}

\textbf{Default --- room as splat (NanoGS in UE 5.8).}  For an indoor room the
recommended default is to represent the room as a real Gaussian splat rendered
by the NanoGS UE 5.8 plugin, with object FBXs from TRELLIS.2 embedded
alongside it.  This is simultaneously the highest-fidelity path and the one
that leans most directly on the vendored LichtFeld trainer: the native
\code{splat\_30000.ply} is cleaned with SuperSplat, imported into UE 5.8, and
rendered as real Gaussians.  The mandated polygonal mesh is satisfied by the
TRELLIS.2 object FBXs.

\textbf{When to use a polygonal room mesh.}  If the brief requires a fully
polygonal game-asset scene (Nanite room mesh, not a splat), use \textbf{2DGS}
(quick win, same Open3D-TSDF backend as the current fallback, but 2D-surfel
depth is much cleaner than ellipsoid depth on flat walls) or \textbf{PGSR}
(planar Gaussians, explicitly handles textureless walls, strongest for indoor
rooms).  The legacy CoMe/Open3D-TSDF combination is retained as a fallback
but produces visibly blockier results and should not be the default.

\textbf{When ArtiFixer helps.}  ArtiFixer is a 3DGS refiner targeting
floaters, holes, and under-observed regions: it applies an opacity-mix
enhancement that fills in areas seen from too few viewpoints.  It runs on a
single 48\,GB RTX 6000 Ada (peak $\approx$\,45.5\,GB) and is the correct tool
for \emph{sharp but coverage-limited} captures.  It does \emph{not} deblur.
Applying ArtiFixer to a blur-limited capture wastes GPU time without improving
results.

\textbf{NanoGS floaters.}  The NanoGS plugin renders the splat faithfully,
including any floater Gaussians in the source PLY.  Pre-clean the PLY with
SuperSplat (floater removal by opacity and scale thresholds) before import.
If floaters persist after SuperSplat, check whether coverage was the issue
(ArtiFixer) or whether the source footage was simply blurry (recapture).

\subsection{Object reconstruction}
\label{sec:pb-objects}

Objects follow the same path regardless of room profile:

\begin{enumerate}
  \item \textbf{SAM3 concept segmentation} identifies freestanding objects in
        the splat and extracts per-object Gaussian PLY subsets.
  \item \textbf{Orbit render}: four synthetic camera views around the object
        centroid are rendered from the Gaussian subset.
  \item \textbf{TRELLIS.2 \code{hull\_e2e}}: the orbit renders are submitted to
        TRELLIS.2 (via ComfyUI at \code{http://vitrine-comfyui:8188}) which runs
        FLUX.2 turnaround generation followed by TRELLIS.2 3D reconstruction,
        producing a textured GLB with PBR materials.
  \item \textbf{Prescale}: Blender prescales the GLB to physical units to match
        the COLMAP coordinate frame.
  \item \textbf{FBX export}: the prescaled mesh is exported as FBX for UE 5.8
        game-asset import (Nanite-ready).
\end{enumerate}

\noindent TRELLIS.2 is the validated primary path; Hunyuan3D-2.1 (also staged
in ComfyUI) is the fallback for cases where the orbit renders are cluttered or
partially occluded.  Restart ComfyUI between objects if VRAM is under pressure:
FLUX.2 and TRELLIS.2 cannot co-reside with other large models on 48\,GB.

\subsection{VRAM phasing}
\label{sec:pb-vram}

The pipeline phases large models to avoid out-of-memory conditions.  The key
constraint is that FLUX.2 and TRELLIS.2 together require $\approx$\,45\,GB; they
cannot share VRAM with other resident models.  The orchestrator manages the
lifecycle automatically (load $\to$ infer $\to$ free) but the operator should
be aware:

\begin{itemize}
  \item If ComfyUI reports an OOM, the previous model was not freed cleanly.
        Restart the \code{vitrine-comfyui} container and re-run from the failing
        object.
  \item Process objects serially, not in parallel, unless a second GPU is
        dedicated to the task.
  \item ArtiFixer peaks at $\approx$\,45.5\,GB; do not run it concurrently with
        any ComfyUI workflow.
\end{itemize}

% ─────────────────────────────────────────────────────────────────────
\section{Troubleshooting}
\label{sec:pb-troubleshoot}
% ─────────────────────────────────────────────────────────────────────

\subsection{Recapture flag: footage too soft}
\label{sec:pb-ts-recapture}

\begin{technote}
\textbf{Symptom.}  The frame-QA stage returns \code{VideoQualityVerdict.too\_low\_quality} for all
or most input videos.  The MUSIQ median is below $\approx$\,25.
\end{technote}

The QA gate cannot manufacture sharpness that was never recorded.  Lowering
the MUSIQ threshold feeds more frames to COLMAP but does not improve the
reconstruction --- it just postpones a failure that COLMAP will surface as low
registration or that 3DGS training will surface as a blurry, noisy splat.

\textbf{Resolution: recapture} following the camera settings and movement
guidance in \S\ref{sec:pb-camera-settings}--\ref{sec:pb-movement}.  Checklist:
\begin{itemize}[nosep]
  \item Switch to 60\,fps (halves exposure time vs.\ 30\,fps).
  \item Lock shutter at $\geq$\,1/500\,s in a pro camera app.
  \item Reduce movement speed; add micro-pauses at key viewpoints.
  \item Use a gimbal or brace the camera firmly.
  \item Ensure bright, even illumination (adequate light allows fast shutter
        without under-exposure).
\end{itemize}

\subsection{Directional blur gate}
\label{sec:pb-ts-blur-gate}

Vitrine applies a directional blur metric in addition to the MUSIQ score:
optical-flow magnitude and direction identify frames where camera motion was
predominantly linear (walking forward) versus rotational (panning).  Panning
motion at low shutter speeds produces pronounced horizontal or vertical blur
streaks that fool the Laplacian sharpness metric (which measures spatial
frequency without direction).  Frames with high directional blur are downgraded
before FIFO windowed selection.

If the frame-QA gate discards a high fraction of frames despite the footage
appearing ``mostly sharp'', check whether long pan segments are responsible:
inspect the flow field overlays saved to \code{output/JOB\_ID/qa\_debug/}.
Recapture with slower, more deliberate panning ($<$\,15°/s) or replace panning
segments with lateral translation.

\subsection{COLMAP: low registration rate}
\label{sec:pb-ts-colmap}

A COLMAP registration below 70\,\% usually indicates one of the following:

\begin{description}
  \item[Insufficient overlap] You moved too fast or the scene was only captured
    from one direction.  Recapture with slower movement and the systematic
    multi-pass patterns described in \S\ref{sec:pb-patterns}.
  \item[COLMAP invocation flags] With COLMAP 4.1.0 the option namespace changed;
    old recipes silently use wrong flags.  Use \code{--FeatureExtraction.*} and
    \code{--FeatureMatching.*} (not the old \code{--SiftExtraction.*} namespace).
    Run feature extraction on GPU (\code{--FeatureExtraction.type SIFT} on GPU):
    GPU SIFT achieved $\approx$\,100\,\% registration on the multi-sweep
    dreamlab captures.
  \item[Missing output directory] COLMAP silently fails to create its database if
    the output directory does not already exist.  Run
    \code{mkdir -p colmap/} before invoking \code{feature\_extractor}.
  \item[Mixed cameras] If the capture combined multiple devices, pass
    \code{--ImageReader.single\_camera\_per\_folder} so COLMAP does not force
    one shared intrinsic model across devices.
  \item[Featureless surfaces] Large white walls or uniform floors provide no
    feature points.  Add temporary texture to blank surfaces before
    recapturing (\S\ref{sec:pb-scene-prep}).
\end{description}

\subsection{NanoGS floaters in the UE scene}
\label{sec:pb-ts-floaters}

\begin{description}
  \item[Cause: coverage holes.]  Regions seen from fewer than five viewpoints
    generate noisy or un-anchored Gaussians.  Apply ArtiFixer (the correct tool
    for this profile) before running SuperSplat.
  \item[Cause: motion blur in source footage.]  A blurry Gaussian is still a
    Gaussian; NanoGS renders it faithfully as a blurry blob.  ArtiFixer does
    not help here.  The resolution is recapture.
  \item[Prune with SuperSplat before UE import.] SuperSplat's opacity and scale
    thresholds reliably remove stray low-opacity Gaussians that accumulate in
    poorly observed regions.  Run SuperSplat on every PLY before importing into
    UE regardless of whether ArtiFixer was applied.
\end{description}

\subsection{UE import scale and object placement}
\label{sec:pb-ts-ue-scale}

UE 5.8 imports FBX in centimetres by default; the COLMAP/3DGS coordinate
frame is in metres.  Vitrine's \code{blender\_prescale\_objects} step converts
object meshes to physical scale before FBX export.  If imported objects appear
100$\times$ too large or too small, check that \code{blender\_prescale\_objects}
ran successfully and that the UE import dialogue has \textit{Import Uniform
Scale} set to 1.0 (not 0.01 or 100).  Use \code{ue\_place\_prescaled} (the
keeper placement script) rather than hand-placing objects via Remote Control,
which has been observed to produce incorrect scale on the final assembly step.

\subsection{GPU VRAM OOM during ComfyUI object processing}
\label{sec:pb-ts-vram}

FLUX.2 and TRELLIS.2 together peak at $\approx$\,45.5\,GB.  If a workflow
fails with an out-of-memory error:
\begin{enumerate}[nosep]
  \item Restart the \code{vitrine-comfyui} container to flush all resident
        models: \code{docker restart vitrine-comfyui}.
  \item Re-submit the failing object workflow.  ComfyUI's queue is persistent
        across restarts; the job will resume from the failed node.
  \item If OOM recurs, check that ArtiFixer is not concurrently resident.
        ArtiFixer should be run as a separate stage with ComfyUI idle.
  \item If the \code{vitrine-comfyui} container is sharing GPU\,0 with
        \code{gaussian-toolkit} 3DGS training, defer object processing until
        training is complete.
\end{enumerate}

\noindent The \code{--local\_attn\_size} flag in ArtiFixer is the dominant
VRAM lever: reducing it from the default to a smaller window size (e.g.\ 128)
significantly reduces peak usage at a modest quality cost.  See
\code{docker/artifixer/patches/} for the confirmed fix applied to the
dreamlab run.

% ─────────────────────────────────────────────────────────────────────
\section{Pre-Capture Checklist}
\label{sec:pb-checklist}
% ─────────────────────────────────────────────────────────────────────

Print and check before every capture session.

\begin{center}
\fbox{\begin{minipage}{0.92\textwidth}
\textbf{Pre-capture}\\[2pt]
$\square$\ Scene is fully static: no people, pets, fans, active screens, or flowing water\\
$\square$\ Lighting is even and consistent (meter swing $<$\,1.5 stops over 360°)\\
$\square$\ Mirrors and reflective surfaces covered or removed\\
$\square$\ Featureless walls have temporary texture added\\
$\square$\ Camera battery charged; storage space verified\\[6pt]
\textbf{Camera settings}\\[2pt]
$\square$\ Resolution: 4K; frame rate: 60\,fps\\
$\square$\ Focus: MANUAL, mid-distance\\
$\square$\ Exposure: MANUAL, locked for average scene brightness\\
$\square$\ Shutter: $\geq$\,1/500\,s (locked in pro app)\\
$\square$\ White balance: MANUAL preset matched to light source\\
$\square$\ Lens: 24--50\,mm equivalent (1$\times$ main camera on phones)\\[6pt]
\textbf{Capture}\\[2pt]
$\square$\ Using gimbal or braced two-handed stabilisation\\
$\square$\ Movement speed: $\leq$\,0.5\,m/s; micro-pauses at key viewpoints\\
$\square$\ Perimeter passes at three heights (low / eye / high)\\
$\square$\ Centre diagonal cross-passes\\
$\square$\ Dedicated 360° orbit of each key object at three heights\\
$\square$\ Detail pass of important features (30--50\,cm, 10--15\,s each)\\
$\square$\ Total duration: 60--120\,s per room; 30--60\,s per object\\[6pt]
\textbf{Post-capture}\\[2pt]
$\square$\ Transfer original file via USB or AirDrop (not cloud)\\
$\square$\ No re-encoding unless unavoidable; if required: H.264 CRF\,18, original resolution\\
$\square$\ File placed in \code{data/raw/} on the pipeline host\\
\end{minipage}}
\end{center}

%
% Dr John O'Hare, DreamLab AI Consulting Ltd — June 2026

% =====================================================================
\chapter{Operator Playbook: Running Vitrine End-to-End}
\label{app:playbook}
% =====================================================================

This appendix is a practical guide for a new operator.  It assumes the Docker
stack is accessible but does not assume familiarity with the internals; the
steps call the pipeline through its published Python and REST interfaces.
Read the capture section (\S\ref{sec:pb-capture}) before anything else:
\emph{capture quality is the single largest determinant of output quality},
and no downstream step can compensate for footage that was never sharp.

% ─────────────────────────────────────────────────────────────────────
\section{Capture Methodology: The Dominant Quality Lever}
\label{sec:pb-capture}
% ─────────────────────────────────────────────────────────────────────

\begin{keyfinding}
Motion blur is the dominant failure mode in practice.  Empirical measurement
on the \emph{dreamlab} dataset found zero under-observed area (median
$\approx$\,90 cameras per surface point) yet a holey, noisy reconstruction.
The problem was dense-but-motion-blurred footage, not insufficient coverage.
Blur cannot be recovered in post: selection can raise the sharpness floor, but
it cannot create detail the sensor never recorded.
\end{keyfinding}

\subsection{Camera settings}
\label{sec:pb-camera-settings}

Lock all automatic adjustments before pressing record.  Any automatic change
between frames is interpreted by COLMAP as a change in scene appearance and
degrades feature matching.

\begin{center}
\begin{tabular}{@{}L{3.2cm}L{4.8cm}L{5.8cm}@{}}
\toprule
\textbf{Setting} & \textbf{Recommendation} & \textbf{Why} \\
\midrule
Resolution     & 4K (3840$\times$2160)           & More features per frame; 8K adds storage cost without proportional gain as training downsamples anyway \\
Frame rate     & \textbf{60\,fps}                 & Caps exposure to $\leq$\,1/60\,s vs.\ 1/30\,s at 30\,fps, halving motion blur; also doubles the candidate frames for sharpness selection \\
Focus          & Manual, mid-distance (3--5\,m)  & Autofocus hunting softens frames and breaks feature matching \\
Exposure       & Manual, locked                  & Auto-exposure causes frame-to-frame brightness ramps that COLMAP interprets as surface-appearance changes \\
White balance  & Manual preset                   & Auto white balance shifts colour temperature between frames \\
Shutter speed  & $\geq$\,1/500\,s (lock in pro app) & Short exposure is the primary anti-blur lever; bright, even light lets you achieve this without under-exposure \\
Lens           & 24--50\,mm equivalent (1$\times$ main camera on phones) & Minimal distortion; ultra-wide ($<$16\,mm) fails in periphery; telephoto reduces parallax \\
\bottomrule
\end{tabular}
\end{center}

\medskip
\noindent\textbf{Phone-specific note.}  Flagship phones (iPhone 15 Pro,
Pixel 9 Pro, Samsung S24 Ultra) produce usable video \emph{only} if you use
the main 1$\times$ camera, lock exposure, focus, and white balance via a pro
app (e.g.\ mcpro24fps, FiLMiC Pro), and shoot in good light.  The Pixel 9 Pro
experiment found \textbf{4K/60\,fps} to be the best video mode: 60\,fps caps
exposure and doubles candidate frames.  \textbf{8K/30\,fps is worse}: the
longer exposure from 30\,fps produces blurrier frames; extra pixels on a
blurred frame are wasted.  If the camera supports RAW photo bursts, those are
even better --- very short exposure, no inter-frame compression.

\subsection{Movement and stabilisation}
\label{sec:pb-movement}

\begin{itemize}
  \item \textbf{Maximum walking speed: 0.5\,m/s.}  Counting ``one-Mississippi''
        per step is a practical governor.  When in doubt, slow down further.
  \item \textbf{Use a gimbal} (DJI RS, Zhiyun) or a tripod dolly.  These
        eliminate micro-shake and rolling-shutter artefacts.  Two-handed
        handheld capture is marginal in bright light; single-hand phone capture
        introduces motion blur on most frames and should be avoided.
  \item \textbf{Insert micro-pauses} at key viewpoints.  The sharpness
        sampler harvests the sharp still frames from these pauses.  Even a
        fraction of a second of steadiness at a corner or interesting object is
        worth more than continuous motion.
  \item \textbf{Move laterally, not just forward.}  The pipeline needs
        triangulation, and triangulation needs physical baseline between
        viewpoints.  Sidestep, arc, and weave rather than walking straight ahead.
  \item \textbf{Every point in the scene should be visible from at least five
        viewpoints.}  Revisit areas from different angles rather than passing
        through once.
\end{itemize}

\subsection{Capture patterns}
\label{sec:pb-patterns}

\paragraph{Room scan (60--120\,s per room).}
Walk the full perimeter at three height bands — low ($\approx$\,1\,m, waist),
eye level ($\approx$\,1.6\,m), and high (arms raised, $\approx$\,2.2\,m) —
then walk the two diagonals to fill the interior.  This gives three elevation
bands of every wall plus cross-coverage of floor and ceiling.  Multiple
independent sweeps of the same room at different heights combine into a denser
reconstruction than any single sweep: an observed three-sweep capture produced
$\approx$146\,k COLMAP points at $\approx$95\,\% registration versus $\approx$90\,k
at $\approx$90\,\% for a single perimeter pass.

\paragraph{Object orbit (30--60\,s per object).}
The best room sweep is rarely the best for any individual object.  Objects you
want as individual hulls need a \emph{dedicated} close orbit in addition to
the room sweeps:
\begin{enumerate}[nosep]
  \item Orbit at eye level ($\approx$1.5\,m distance): full 360°, 20\,s.
  \item Orbit at knee height, camera angled upward: full 360°, 20\,s.
  \item Orbit above head, camera angled downward: full 360°, 20\,s.
\end{enumerate}

\paragraph{Detail pass (10--15\,s per feature).}
After room and object passes, sweep 30--50\,cm from important surfaces (carved
detail, lettering, fine texture).  These close-up frames lift the resolution
ceiling for that region.

\subsection{Scene preparation}
\label{sec:pb-scene-prep}

\begin{itemize}
  \item \textbf{Remove all moving elements.}  People, animals, ceiling fans,
        screens playing video, swaying plants — any inter-frame motion creates
        smeared ghost geometry.  Pause recording if someone walks through the
        scene.  Even a person standing still shifts weight and breathes.
  \item \textbf{Turn screens off.}  A powered-off TV is far less glossy than a
        running one.  A running TV is the worst-case object: a large,
        view-dependent reflector whose reflections defeat both SfM (inconsistent
        features) and 3DGS (spherical-harmonics overfit the moving reflection
        into floaters).  The 2026-06-21 Drive dataset was abandoned for exactly
        this reason combined with low MUSIQ scores.
  \item \textbf{Cover mirrors and polished surfaces.}  Matte fabric or dulling
        spray prevents ghost geometry behind reflective surfaces.
  \item \textbf{Ensure even, diffuse lighting.}  Overcast sky is ideal outdoors.
        Indoors, use large soft panels or bounce lights off the ceiling.  Before
        recording, rotate 360° in the scene centre watching the exposure meter:
        if it swings more than 1.5 stops, add fill light or flag bright sources.
  \item \textbf{Add temporary texture to featureless surfaces.}  White walls and
        plain ceilings give COLMAP nothing to match.  Place books, plants, or
        patterned cloth against blank walls during capture; remove them later if
        needed.
\end{itemize}

\subsection{Do / do not summary}
\label{sec:pb-donot}

\begin{center}
\begin{tabular}{@{}L{5.8cm}L{7.2cm}@{}}
\toprule
\textbf{Do} & \textbf{Do not} \\
\midrule
4K/60\,fps, manual exposure, focus, WB & Use auto-exposure or auto white balance \\
Gimbal or two-handed stabilisation & Pan fast ($>$15°/s) or shoot single-handed \\
$\leq$\,0.5\,m/s, micro-pauses at key points & Walk at normal pace \\
Three height bands + diagonals & Walk only one straight path through the space \\
Dedicated 360° object orbits & Assume the room sweep covers individual objects \\
Transfer the original camera file over USB & Upload to a cloud service that transcodes \\
Even, diffuse lighting & Mix tungsten and daylight; allow hard shadows \\
Turn screens and fans off & Capture with any motion in the scene \\
$\geq$1/500\,s shutter (lock in pro app) & Rely on in-camera optical stabilisation alone for blur \\
4K/60\,fps on phones & Use 8K/30\,fps (blurrier despite higher pixel count) \\
\bottomrule
\end{tabular}
\end{center}

\subsection{Why blur is unrecoverable}
\label{sec:pb-blur-unrecoverable}

Motion blur is a measurement artefact: a blurred pixel is the average of many
spatial samples acquired over the exposure window, and the individual samples
are gone.  Per-frame generative deblurring (NAFNet, Restormer) hallucinates
texture rather than recovering it, and the hallucinated texture is spatially
inconsistent across frames, which breaks multi-view SfM.  The only
reconstruction methods that legitimately model blur --- joint exposure-trajectory
optimisation approaches such as \textsc{3dgs-deblur} --- show gains of
approximately 1--2\,dB PSNR on \emph{moderate} blur but hit the same
SfM-quality ceiling as the source footage.  \textbf{Blur is unrecoverable in
post; the real lever is short-shutter capture.}

The Vitrine frame-QA gate applies a MUSIQ no-reference IQA score plus a full-resolution
variance-of-Laplacian sharpness score to every extracted frame.  A video that cannot
supply \code{min\_good\_frames} above the KEEP threshold is marked
\code{too\_low\_quality} and the operator is asked to recapture.  The dreamlab
dataset scored a median MUSIQ of $\approx$\,31 (usable but blur-limited); the
abandoned Drive dataset scored $\approx$\,19 (unusable).  The gate is a diagnostic,
not a repair: it tells you when recapture is necessary.

% ─────────────────────────────────────────────────────────────────────
\section{Running the Pipeline}
\label{sec:pb-run}
% ─────────────────────────────────────────────────────────────────────

\subsection{Bringing up the stack}
\label{sec:pb-stack}

All pipeline work runs inside the \code{gaussian-toolkit} container.  From the
repository root on the host:

\begin{lstlisting}[language=bash,caption={Bring up the consolidated stack}]
# Clone with the vendored LichtFeld tool included
git clone --recurse-submodules <repo-url>
cd Vitrine

# Start the pipeline containers (GPU required)
docker compose -f docker-compose.consolidated.yml up -d

# Verify the stack is healthy
curl http://gaussian-toolkit:7860/health     # web UI
curl http://gaussian-toolkit:45677/mcp       # LichtFeld MCP
curl http://vitrine-comfyui:8188/system_stats  # ComfyUI (TRELLIS.2 / FLUX.2)
\end{lstlisting}

\noindent The Unreal Engine 5.8 overlay is optional and started separately:

\begin{lstlisting}[language=bash,caption={Bring up the Unreal overlay (optional)}]
docker compose \
  -f docker-compose.consolidated.yml \
  -f unreal/docker-compose.unreal.yml \
  up -d unreal unreal-mcp-bridge
\end{lstlisting}

\noindent Run the SOTA preflight check to verify model weights and VRAM before
processing any footage:

\begin{lstlisting}[language=bash,caption={SOTA preflight}]
python -m pipeline.sota_registry check
\end{lstlisting}

\noindent If any weight or pin is stale the preflight reports it; resolve before
continuing.  Model weights live in \code{data/comfyui/models/} ($\approx$\,216\,GB,
bind-mounted read-only); do not re-download them.

\subsection{Ingest and frame QA}
\label{sec:pb-ingest}

Place source video files in \code{data/raw/} and submit the job via the web UI
(\texttt{http://localhost:7860}) or directly:

\begin{lstlisting}[language=bash,caption={Step 1 -- ingest and frame extraction}]
python3 -c "
from pipeline.stages import PipelineStages
p = PipelineStages('/data/output/JOB_ID')
result = p.ingest('/data/raw/scene.mp4', fps=4.0)
print(result)
"
\end{lstlisting}

\noindent Ingest extracts frames at 4\,fps to oversample, then the frame-QA
stage selects the sharpest subset.  The QA pipeline applies two complementary
metrics:

\begin{enumerate}
  \item \textbf{MUSIQ no-reference IQA} (\code{pyiqa}, GPU).  A
        perceptual score on the KonIQ-10k scale ($\approx$0--100; good
        photography scores 50--75; severely blurred or noisy footage scores
        below $\approx$25).  Videos unable to supply enough frames above the
        KEEP threshold are marked \code{too\_low\_quality} and skipped.
  \item \textbf{Full-resolution variance-of-Laplacian.}  Sharpness is measured
        at full 4K resolution, not at a downscaled proxy --- downscaling blurs
        away the motion-blur signal (960p blurdetect showed only a
        $\approx$1.6$\times$ range on dreamlab frames; full-4K Laplacian showed a
        $\approx$105$\times$ range).  The \emph{sharpest frame per overlapping
        FIFO window} is kept, which raises the sharpness floor by 20--38\,\% at
        the bottom quartile --- the frames that most poison COLMAP matches.
\end{enumerate}

\noindent Inspect the selected frame count before proceeding:

\begin{lstlisting}[language=bash,caption={Step 3 -- select best frames}]
python3 -c "
from pipeline.stages import PipelineStages
p = PipelineStages('/data/output/JOB_ID')
result = p.select_frames('/data/output/JOB_ID/frames/', target=80)
print(result)
"
# Target: 60-80 frames. If < 40 are above threshold, consider recapture.
\end{lstlisting}

\subsection{Structure-from-Motion and 3DGS training}
\label{sec:pb-sfm-train}

\begin{lstlisting}[language=bash,caption={Step 4 -- COLMAP SfM (sequential matcher for video)}]
python3 -c "
from pipeline.stages import PipelineStages
p = PipelineStages('/data/output/JOB_ID')
result = p.reconstruct('/data/output/JOB_ID/frames_selected/',
                       matcher='sequential')
print(result)
"
# Quality gate: >= 70 % of frames must register.
# If < 50 % register, re-run with matcher='exhaustive'.
\end{lstlisting}

\begin{lstlisting}[language=bash,caption={Step 5 -- 3DGS training (30\,000 iterations)}]
python3 -c "
from pipeline.stages import PipelineStages
p = PipelineStages('/data/output/JOB_ID')
result = p.train('/data/output/JOB_ID/colmap/undistorted/',
                 iterations=30000)
print(result)
"
# Quality gate: PLY file should be > 10 MB; loss < 0.02.
\end{lstlisting}

\noindent The output splat is \code{output/JOB\_ID/model/splat\_30000.ply}.  This
file is the primary LichtFeld-native deliverable that feeds both the Gaussian
splat room path and the object segmentation branch.

\subsection{The stage sequence}
\label{sec:pb-sequence}

The full pipeline in order:

\begin{center}
\begin{tabular}{@{}lL{4.5cm}L{6cm}@{}}
\toprule
\textbf{Stage} & \textbf{What runs} & \textbf{Quality check} \\
\midrule
1 & Ingest: frame extraction at 4\,fps & \code{frames/} directory populated \\
2 & Remove people (if present) & Inspect sample frames \\
3 & Frame-QA selection & 60--80 KEEP-grade frames \\
4 & COLMAP SfM (sequential matcher) & $\geq$\,70\,\% registration \\
5 & 3DGS training (30\,k iterations) & PLY $>$10\,MB, loss $<$0.02 \\
6 & SAM3 concept segmentation & $\geq$\,2 objects identified \\
7 & Per-object PLY extraction + orbit renders & $\geq$5\,k vertices per object \\
7c & TRELLIS.2 hull\_e2e (textured GLB) & Textured GLB produced per object \\
8 & Room mesh extraction (TSDF or 2DGS) & $\geq$\,30\,k vertices, clean topology \\
8b & Mesh cleanup (\code{mesh\_cleaner smooth=0}) & No stray bridging triangles \\
8c & Texture bake (xatlas UV + vertex colours) & Texture atlas per mesh \\
9 & FBX export (\code{blender\_obj\_to\_fbx}) & FBX files produced \\
9b & UE 5.8 import (Nanite game assets) & Assets visible in Content Browser \\
10 & Validate + job complete & Full scene: room + objects \\
\bottomrule
\end{tabular}
\end{center}

\noindent \textbf{Do not stop after training.}  The pipeline is only complete
when it reaches the textured UE scene; stopping at the splat leaves the
deliverable unfinished.

% ─────────────────────────────────────────────────────────────────────
\section{Per-Capture Recipe Selection}
\label{sec:pb-recipe}
% ─────────────────────────────────────────────────────────────────────

Vitrine is not a fixed pipeline.  Each capture is \emph{diagnosed} first; the
router selects the configuration that best fits the capture's bottleneck.  The
four diagnostic axes are:

\begin{description}
  \item[Sharpness] Measured by MUSIQ + full-resolution Laplacian.  Is the
    footage blur-limited?
  \item[Coverage] Measured by a per-surface-point viewpoint count.  Are there
    under-observed regions (holes) or is coverage uniformly dense?
  \item[Hardware] RGB camera only, depth sensor (iPhone LiDAR), or event camera?
  \item[Deliverable] Fidelity-first splat room, polygonal mesh room, or both?
\end{description}

\noindent Table~\ref{tab:pb-recipes} maps common capture profiles to their
recommended configurations.

\begin{table}[ht]
\centering
\caption{Capture profile to recipe selector.  Objects are always reconstructed
via TRELLIS.2 \code{hull\_e2e} regardless of room profile.}
\label{tab:pb-recipes}
\begin{tabular}{@{}L{3.4cm}L{2.4cm}L{2.4cm}L{2.2cm}L{3.0cm}@{}}
\toprule
\textbf{Capture profile} & \textbf{Room rep} & \textbf{Enhancer} & \textbf{Mesh backend} & \textbf{Note} \\
\midrule
\textbf{Blurry $\cdot$ dense $\cdot$ RGB}
  (e.g.\ dreamlab)
  & Splat hero + mesh asset (dual)
  & \emph{None} (ArtiFixer won't help)
  & 2DGS / PGSR
  & Recapture is the real ceiling \\
\addlinespace
\textbf{Sharp $\cdot$ holey $\cdot$ RGB}
  & Splat (primary)
  & \textbf{ArtiFixer} (fills floaters/holes)
  & 2DGS / PGSR if mesh required
  & ArtiFixer's sweet spot \\
\addlinespace
\textbf{Sharp $\cdot$ dense $\cdot$ RGB}
  & Splat (primary)
  & None
  & 2DGS if mesh needed
  & Best case; least pain \\
\addlinespace
\textbf{iPhone LiDAR depth}
  & Mesh (primary)
  & Optional
  & \textbf{AGS-Mesh} (sensor depth)
  & Depth sensor unlocks clean walls \\
\addlinespace
\textbf{Event camera $\cdot$ blurry}
  & Splat
  & ---
  & ---
  & EBAD-Gaussian deblurs at source \\
\bottomrule
\end{tabular}
\end{table}

\subsection{Splat vs.\ mesh for the room}
\label{sec:pb-splat-vs-mesh}

\textbf{Default --- room as splat (NanoGS in UE 5.8).}  For an indoor room the
recommended default is to represent the room as a real Gaussian splat rendered
by the NanoGS UE 5.8 plugin, with object FBXs from TRELLIS.2 embedded
alongside it.  This is simultaneously the highest-fidelity path and the one
that leans most directly on the vendored LichtFeld trainer: the native
\code{splat\_30000.ply} is cleaned with SuperSplat, imported into UE 5.8, and
rendered as real Gaussians.  The mandated polygonal mesh is satisfied by the
TRELLIS.2 object FBXs.

\textbf{When to use a polygonal room mesh.}  If the brief requires a fully
polygonal game-asset scene (Nanite room mesh, not a splat), use \textbf{2DGS}
(quick win, same Open3D-TSDF backend as the current fallback, but 2D-surfel
depth is much cleaner than ellipsoid depth on flat walls) or \textbf{PGSR}
(planar Gaussians, explicitly handles textureless walls, strongest for indoor
rooms).  The legacy CoMe/Open3D-TSDF combination is retained as a fallback
but produces visibly blockier results and should not be the default.

\textbf{When ArtiFixer helps.}  ArtiFixer is a 3DGS refiner targeting
floaters, holes, and under-observed regions: it applies an opacity-mix
enhancement that fills in areas seen from too few viewpoints.  It runs on a
single 48\,GB RTX 6000 Ada (peak $\approx$\,45.5\,GB) and is the correct tool
for \emph{sharp but coverage-limited} captures.  It does \emph{not} deblur.
Applying ArtiFixer to a blur-limited capture wastes GPU time without improving
results.

\textbf{NanoGS floaters.}  The NanoGS plugin renders the splat faithfully,
including any floater Gaussians in the source PLY.  Pre-clean the PLY with
SuperSplat (floater removal by opacity and scale thresholds) before import.
If floaters persist after SuperSplat, check whether coverage was the issue
(ArtiFixer) or whether the source footage was simply blurry (recapture).

\subsection{Object reconstruction}
\label{sec:pb-objects}

Objects follow the same path regardless of room profile:

\begin{enumerate}
  \item \textbf{SAM3 concept segmentation} identifies freestanding objects in
        the splat and extracts per-object Gaussian PLY subsets.
  \item \textbf{Orbit render}: four synthetic camera views around the object
        centroid are rendered from the Gaussian subset.
  \item \textbf{TRELLIS.2 \code{hull\_e2e}}: the orbit renders are submitted to
        TRELLIS.2 (via ComfyUI at \code{http://vitrine-comfyui:8188}) which runs
        FLUX.2 turnaround generation followed by TRELLIS.2 3D reconstruction,
        producing a textured GLB with PBR materials.
  \item \textbf{Prescale}: Blender prescales the GLB to physical units to match
        the COLMAP coordinate frame.
  \item \textbf{FBX export}: the prescaled mesh is exported as FBX for UE 5.8
        game-asset import (Nanite-ready).
\end{enumerate}

\noindent TRELLIS.2 is the validated primary path; Hunyuan3D-2.1 (also staged
in ComfyUI) is the fallback for cases where the orbit renders are cluttered or
partially occluded.  Restart ComfyUI between objects if VRAM is under pressure:
FLUX.2 and TRELLIS.2 cannot co-reside with other large models on 48\,GB.

\subsection{VRAM phasing}
\label{sec:pb-vram}

The pipeline phases large models to avoid out-of-memory conditions.  The key
constraint is that FLUX.2 and TRELLIS.2 together require $\approx$\,45\,GB; they
cannot share VRAM with other resident models.  The orchestrator manages the
lifecycle automatically (load $\to$ infer $\to$ free) but the operator should
be aware:

\begin{itemize}
  \item If ComfyUI reports an OOM, the previous model was not freed cleanly.
        Restart the \code{vitrine-comfyui} container and re-run from the failing
        object.
  \item Process objects serially, not in parallel, unless a second GPU is
        dedicated to the task.
  \item ArtiFixer peaks at $\approx$\,45.5\,GB; do not run it concurrently with
        any ComfyUI workflow.
\end{itemize}

% ─────────────────────────────────────────────────────────────────────
\section{Troubleshooting}
\label{sec:pb-troubleshoot}
% ─────────────────────────────────────────────────────────────────────

\subsection{Recapture flag: footage too soft}
\label{sec:pb-ts-recapture}

\begin{technote}
\textbf{Symptom.}  The frame-QA stage returns \code{VideoQualityVerdict.too\_low\_quality} for all
or most input videos.  The MUSIQ median is below $\approx$\,25.
\end{technote}

The QA gate cannot manufacture sharpness that was never recorded.  Lowering
the MUSIQ threshold feeds more frames to COLMAP but does not improve the
reconstruction --- it just postpones a failure that COLMAP will surface as low
registration or that 3DGS training will surface as a blurry, noisy splat.

\textbf{Resolution: recapture} following the camera settings and movement
guidance in \S\ref{sec:pb-camera-settings}--\ref{sec:pb-movement}.  Checklist:
\begin{itemize}[nosep]
  \item Switch to 60\,fps (halves exposure time vs.\ 30\,fps).
  \item Lock shutter at $\geq$\,1/500\,s in a pro camera app.
  \item Reduce movement speed; add micro-pauses at key viewpoints.
  \item Use a gimbal or brace the camera firmly.
  \item Ensure bright, even illumination (adequate light allows fast shutter
        without under-exposure).
\end{itemize}

\subsection{Directional blur gate}
\label{sec:pb-ts-blur-gate}

Vitrine applies a directional blur metric in addition to the MUSIQ score:
optical-flow magnitude and direction identify frames where camera motion was
predominantly linear (walking forward) versus rotational (panning).  Panning
motion at low shutter speeds produces pronounced horizontal or vertical blur
streaks that fool the Laplacian sharpness metric (which measures spatial
frequency without direction).  Frames with high directional blur are downgraded
before FIFO windowed selection.

If the frame-QA gate discards a high fraction of frames despite the footage
appearing ``mostly sharp'', check whether long pan segments are responsible:
inspect the flow field overlays saved to \code{output/JOB\_ID/qa\_debug/}.
Recapture with slower, more deliberate panning ($<$\,15°/s) or replace panning
segments with lateral translation.

\subsection{COLMAP: low registration rate}
\label{sec:pb-ts-colmap}

A COLMAP registration below 70\,\% usually indicates one of the following:

\begin{description}
  \item[Insufficient overlap] You moved too fast or the scene was only captured
    from one direction.  Recapture with slower movement and the systematic
    multi-pass patterns described in \S\ref{sec:pb-patterns}.
  \item[COLMAP invocation flags] With COLMAP 4.1.0 the option namespace changed;
    old recipes silently use wrong flags.  Use \code{--FeatureExtraction.*} and
    \code{--FeatureMatching.*} (not the old \code{--SiftExtraction.*} namespace).
    Run feature extraction on GPU (\code{--FeatureExtraction.type SIFT} on GPU):
    GPU SIFT achieved $\approx$\,100\,\% registration on the multi-sweep
    dreamlab captures.
  \item[Missing output directory] COLMAP silently fails to create its database if
    the output directory does not already exist.  Run
    \code{mkdir -p colmap/} before invoking \code{feature\_extractor}.
  \item[Mixed cameras] If the capture combined multiple devices, pass
    \code{--ImageReader.single\_camera\_per\_folder} so COLMAP does not force
    one shared intrinsic model across devices.
  \item[Featureless surfaces] Large white walls or uniform floors provide no
    feature points.  Add temporary texture to blank surfaces before
    recapturing (\S\ref{sec:pb-scene-prep}).
\end{description}

\subsection{NanoGS floaters in the UE scene}
\label{sec:pb-ts-floaters}

\begin{description}
  \item[Cause: coverage holes.]  Regions seen from fewer than five viewpoints
    generate noisy or un-anchored Gaussians.  Apply ArtiFixer (the correct tool
    for this profile) before running SuperSplat.
  \item[Cause: motion blur in source footage.]  A blurry Gaussian is still a
    Gaussian; NanoGS renders it faithfully as a blurry blob.  ArtiFixer does
    not help here.  The resolution is recapture.
  \item[Prune with SuperSplat before UE import.] SuperSplat's opacity and scale
    thresholds reliably remove stray low-opacity Gaussians that accumulate in
    poorly observed regions.  Run SuperSplat on every PLY before importing into
    UE regardless of whether ArtiFixer was applied.
\end{description}

\subsection{UE import scale and object placement}
\label{sec:pb-ts-ue-scale}

UE 5.8 imports FBX in centimetres by default; the COLMAP/3DGS coordinate
frame is in metres.  Vitrine's \code{blender\_prescale\_objects} step converts
object meshes to physical scale before FBX export.  If imported objects appear
100$\times$ too large or too small, check that \code{blender\_prescale\_objects}
ran successfully and that the UE import dialogue has \textit{Import Uniform
Scale} set to 1.0 (not 0.01 or 100).  Use \code{ue\_place\_prescaled} (the
keeper placement script) rather than hand-placing objects via Remote Control,
which has been observed to produce incorrect scale on the final assembly step.

\subsection{GPU VRAM OOM during ComfyUI object processing}
\label{sec:pb-ts-vram}

FLUX.2 and TRELLIS.2 together peak at $\approx$\,45.5\,GB.  If a workflow
fails with an out-of-memory error:
\begin{enumerate}[nosep]
  \item Restart the \code{vitrine-comfyui} container to flush all resident
        models: \code{docker restart vitrine-comfyui}.
  \item Re-submit the failing object workflow.  ComfyUI's queue is persistent
        across restarts; the job will resume from the failed node.
  \item If OOM recurs, check that ArtiFixer is not concurrently resident.
        ArtiFixer should be run as a separate stage with ComfyUI idle.
  \item If the \code{vitrine-comfyui} container is sharing GPU\,0 with
        \code{gaussian-toolkit} 3DGS training, defer object processing until
        training is complete.
\end{enumerate}

\noindent The \code{--local\_attn\_size} flag in ArtiFixer is the dominant
VRAM lever: reducing it from the default to a smaller window size (e.g.\ 128)
significantly reduces peak usage at a modest quality cost.  See
\code{docker/artifixer/patches/} for the confirmed fix applied to the
dreamlab run.

% ─────────────────────────────────────────────────────────────────────
\section{Pre-Capture Checklist}
\label{sec:pb-checklist}
% ─────────────────────────────────────────────────────────────────────

Print and check before every capture session.

\begin{center}
\fbox{\begin{minipage}{0.92\textwidth}
\textbf{Pre-capture}\\[2pt]
$\square$\ Scene is fully static: no people, pets, fans, active screens, or flowing water\\
$\square$\ Lighting is even and consistent (meter swing $<$\,1.5 stops over 360°)\\
$\square$\ Mirrors and reflective surfaces covered or removed\\
$\square$\ Featureless walls have temporary texture added\\
$\square$\ Camera battery charged; storage space verified\\[6pt]
\textbf{Camera settings}\\[2pt]
$\square$\ Resolution: 4K; frame rate: 60\,fps\\
$\square$\ Focus: MANUAL, mid-distance\\
$\square$\ Exposure: MANUAL, locked for average scene brightness\\
$\square$\ Shutter: $\geq$\,1/500\,s (locked in pro app)\\
$\square$\ White balance: MANUAL preset matched to light source\\
$\square$\ Lens: 24--50\,mm equivalent (1$\times$ main camera on phones)\\[6pt]
\textbf{Capture}\\[2pt]
$\square$\ Using gimbal or braced two-handed stabilisation\\
$\square$\ Movement speed: $\leq$\,0.5\,m/s; micro-pauses at key viewpoints\\
$\square$\ Perimeter passes at three heights (low / eye / high)\\
$\square$\ Centre diagonal cross-passes\\
$\square$\ Dedicated 360° orbit of each key object at three heights\\
$\square$\ Detail pass of important features (30--50\,cm, 10--15\,s each)\\
$\square$\ Total duration: 60--120\,s per room; 30--60\,s per object\\[6pt]
\textbf{Post-capture}\\[2pt]
$\square$\ Transfer original file via USB or AirDrop (not cloud)\\
$\square$\ No re-encoding unless unavoidable; if required: H.264 CRF\,18, original resolution\\
$\square$\ File placed in \code{data/raw/} on the pipeline host\\
\end{minipage}}
\end{center}

%
% Dr John O'Hare, DreamLab AI Consulting Ltd — June 2026

% =====================================================================
\chapter{Operator Playbook: Running Vitrine End-to-End}
\label{app:playbook}
% =====================================================================

This appendix is a practical guide for a new operator.  It assumes the Docker
stack is accessible but does not assume familiarity with the internals; the
steps call the pipeline through its published Python and REST interfaces.
Read the capture section (\S\ref{sec:pb-capture}) before anything else:
\emph{capture quality is the single largest determinant of output quality},
and no downstream step can compensate for footage that was never sharp.

% ─────────────────────────────────────────────────────────────────────
\section{Capture Methodology: The Dominant Quality Lever}
\label{sec:pb-capture}
% ─────────────────────────────────────────────────────────────────────

\begin{keyfinding}
Motion blur is the dominant failure mode in practice.  Empirical measurement
on the \emph{dreamlab} dataset found zero under-observed area (median
$\approx$\,90 cameras per surface point) yet a holey, noisy reconstruction.
The problem was dense-but-motion-blurred footage, not insufficient coverage.
Blur cannot be recovered in post: selection can raise the sharpness floor, but
it cannot create detail the sensor never recorded.
\end{keyfinding}

\subsection{Camera settings}
\label{sec:pb-camera-settings}

Lock all automatic adjustments before pressing record.  Any automatic change
between frames is interpreted by COLMAP as a change in scene appearance and
degrades feature matching.

\begin{center}
\begin{tabular}{@{}L{3.2cm}L{4.8cm}L{5.8cm}@{}}
\toprule
\textbf{Setting} & \textbf{Recommendation} & \textbf{Why} \\
\midrule
Resolution     & 4K (3840$\times$2160)           & More features per frame; 8K adds storage cost without proportional gain as training downsamples anyway \\
Frame rate     & \textbf{60\,fps}                 & Caps exposure to $\leq$\,1/60\,s vs.\ 1/30\,s at 30\,fps, halving motion blur; also doubles the candidate frames for sharpness selection \\
Focus          & Manual, mid-distance (3--5\,m)  & Autofocus hunting softens frames and breaks feature matching \\
Exposure       & Manual, locked                  & Auto-exposure causes frame-to-frame brightness ramps that COLMAP interprets as surface-appearance changes \\
White balance  & Manual preset                   & Auto white balance shifts colour temperature between frames \\
Shutter speed  & $\geq$\,1/500\,s (lock in pro app) & Short exposure is the primary anti-blur lever; bright, even light lets you achieve this without under-exposure \\
Lens           & 24--50\,mm equivalent (1$\times$ main camera on phones) & Minimal distortion; ultra-wide ($<$16\,mm) fails in periphery; telephoto reduces parallax \\
\bottomrule
\end{tabular}
\end{center}

\medskip
\noindent\textbf{Phone-specific note.}  Flagship phones (iPhone 15 Pro,
Pixel 9 Pro, Samsung S24 Ultra) produce usable video \emph{only} if you use
the main 1$\times$ camera, lock exposure, focus, and white balance via a pro
app (e.g.\ mcpro24fps, FiLMiC Pro), and shoot in good light.  The Pixel 9 Pro
experiment found \textbf{4K/60\,fps} to be the best video mode: 60\,fps caps
exposure and doubles candidate frames.  \textbf{8K/30\,fps is worse}: the
longer exposure from 30\,fps produces blurrier frames; extra pixels on a
blurred frame are wasted.  If the camera supports RAW photo bursts, those are
even better --- very short exposure, no inter-frame compression.

\subsection{Movement and stabilisation}
\label{sec:pb-movement}

\begin{itemize}
  \item \textbf{Maximum walking speed: 0.5\,m/s.}  Counting ``one-Mississippi''
        per step is a practical governor.  When in doubt, slow down further.
  \item \textbf{Use a gimbal} (DJI RS, Zhiyun) or a tripod dolly.  These
        eliminate micro-shake and rolling-shutter artefacts.  Two-handed
        handheld capture is marginal in bright light; single-hand phone capture
        introduces motion blur on most frames and should be avoided.
  \item \textbf{Insert micro-pauses} at key viewpoints.  The sharpness
        sampler harvests the sharp still frames from these pauses.  Even a
        fraction of a second of steadiness at a corner or interesting object is
        worth more than continuous motion.
  \item \textbf{Move laterally, not just forward.}  The pipeline needs
        triangulation, and triangulation needs physical baseline between
        viewpoints.  Sidestep, arc, and weave rather than walking straight ahead.
  \item \textbf{Every point in the scene should be visible from at least five
        viewpoints.}  Revisit areas from different angles rather than passing
        through once.
\end{itemize}

\subsection{Capture patterns}
\label{sec:pb-patterns}

\paragraph{Room scan (60--120\,s per room).}
Walk the full perimeter at three height bands — low ($\approx$\,1\,m, waist),
eye level ($\approx$\,1.6\,m), and high (arms raised, $\approx$\,2.2\,m) —
then walk the two diagonals to fill the interior.  This gives three elevation
bands of every wall plus cross-coverage of floor and ceiling.  Multiple
independent sweeps of the same room at different heights combine into a denser
reconstruction than any single sweep: an observed three-sweep capture produced
$\approx$146\,k COLMAP points at $\approx$95\,\% registration versus $\approx$90\,k
at $\approx$90\,\% for a single perimeter pass.

\paragraph{Object orbit (30--60\,s per object).}
The best room sweep is rarely the best for any individual object.  Objects you
want as individual hulls need a \emph{dedicated} close orbit in addition to
the room sweeps:
\begin{enumerate}[nosep]
  \item Orbit at eye level ($\approx$1.5\,m distance): full 360°, 20\,s.
  \item Orbit at knee height, camera angled upward: full 360°, 20\,s.
  \item Orbit above head, camera angled downward: full 360°, 20\,s.
\end{enumerate}

\paragraph{Detail pass (10--15\,s per feature).}
After room and object passes, sweep 30--50\,cm from important surfaces (carved
detail, lettering, fine texture).  These close-up frames lift the resolution
ceiling for that region.

\subsection{Scene preparation}
\label{sec:pb-scene-prep}

\begin{itemize}
  \item \textbf{Remove all moving elements.}  People, animals, ceiling fans,
        screens playing video, swaying plants — any inter-frame motion creates
        smeared ghost geometry.  Pause recording if someone walks through the
        scene.  Even a person standing still shifts weight and breathes.
  \item \textbf{Turn screens off.}  A powered-off TV is far less glossy than a
        running one.  A running TV is the worst-case object: a large,
        view-dependent reflector whose reflections defeat both SfM (inconsistent
        features) and 3DGS (spherical-harmonics overfit the moving reflection
        into floaters).  The 2026-06-21 Drive dataset was abandoned for exactly
        this reason combined with low MUSIQ scores.
  \item \textbf{Cover mirrors and polished surfaces.}  Matte fabric or dulling
        spray prevents ghost geometry behind reflective surfaces.
  \item \textbf{Ensure even, diffuse lighting.}  Overcast sky is ideal outdoors.
        Indoors, use large soft panels or bounce lights off the ceiling.  Before
        recording, rotate 360° in the scene centre watching the exposure meter:
        if it swings more than 1.5 stops, add fill light or flag bright sources.
  \item \textbf{Add temporary texture to featureless surfaces.}  White walls and
        plain ceilings give COLMAP nothing to match.  Place books, plants, or
        patterned cloth against blank walls during capture; remove them later if
        needed.
\end{itemize}

\subsection{Do / do not summary}
\label{sec:pb-donot}

\begin{center}
\begin{tabular}{@{}L{5.8cm}L{7.2cm}@{}}
\toprule
\textbf{Do} & \textbf{Do not} \\
\midrule
4K/60\,fps, manual exposure, focus, WB & Use auto-exposure or auto white balance \\
Gimbal or two-handed stabilisation & Pan fast ($>$15°/s) or shoot single-handed \\
$\leq$\,0.5\,m/s, micro-pauses at key points & Walk at normal pace \\
Three height bands + diagonals & Walk only one straight path through the space \\
Dedicated 360° object orbits & Assume the room sweep covers individual objects \\
Transfer the original camera file over USB & Upload to a cloud service that transcodes \\
Even, diffuse lighting & Mix tungsten and daylight; allow hard shadows \\
Turn screens and fans off & Capture with any motion in the scene \\
$\geq$1/500\,s shutter (lock in pro app) & Rely on in-camera optical stabilisation alone for blur \\
4K/60\,fps on phones & Use 8K/30\,fps (blurrier despite higher pixel count) \\
\bottomrule
\end{tabular}
\end{center}

\subsection{Why blur is unrecoverable}
\label{sec:pb-blur-unrecoverable}

Motion blur is a measurement artefact: a blurred pixel is the average of many
spatial samples acquired over the exposure window, and the individual samples
are gone.  Per-frame generative deblurring (NAFNet, Restormer) hallucinates
texture rather than recovering it, and the hallucinated texture is spatially
inconsistent across frames, which breaks multi-view SfM.  The only
reconstruction methods that legitimately model blur --- joint exposure-trajectory
optimisation approaches such as \textsc{3dgs-deblur} --- show gains of
approximately 1--2\,dB PSNR on \emph{moderate} blur but hit the same
SfM-quality ceiling as the source footage.  \textbf{Blur is unrecoverable in
post; the real lever is short-shutter capture.}

The Vitrine frame-QA gate applies a MUSIQ no-reference IQA score plus a full-resolution
variance-of-Laplacian sharpness score to every extracted frame.  A video that cannot
supply \code{min\_good\_frames} above the KEEP threshold is marked
\code{too\_low\_quality} and the operator is asked to recapture.  The dreamlab
dataset scored a median MUSIQ of $\approx$\,31 (usable but blur-limited); the
abandoned Drive dataset scored $\approx$\,19 (unusable).  The gate is a diagnostic,
not a repair: it tells you when recapture is necessary.

% ─────────────────────────────────────────────────────────────────────
\section{Running the Pipeline}
\label{sec:pb-run}
% ─────────────────────────────────────────────────────────────────────

\subsection{Bringing up the stack}
\label{sec:pb-stack}

All pipeline work runs inside the \code{gaussian-toolkit} container.  From the
repository root on the host:

\begin{lstlisting}[language=bash,caption={Bring up the consolidated stack}]
# Clone with the vendored LichtFeld tool included
git clone --recurse-submodules <repo-url>
cd Vitrine

# Start the pipeline containers (GPU required)
docker compose -f docker-compose.consolidated.yml up -d

# Verify the stack is healthy
curl http://gaussian-toolkit:7860/health     # web UI
curl http://gaussian-toolkit:45677/mcp       # LichtFeld MCP
curl http://vitrine-comfyui:8188/system_stats  # ComfyUI (TRELLIS.2 / FLUX.2)
\end{lstlisting}

\noindent The Unreal Engine 5.8 overlay is optional and started separately:

\begin{lstlisting}[language=bash,caption={Bring up the Unreal overlay (optional)}]
docker compose \
  -f docker-compose.consolidated.yml \
  -f unreal/docker-compose.unreal.yml \
  up -d unreal unreal-mcp-bridge
\end{lstlisting}

\noindent Run the SOTA preflight check to verify model weights and VRAM before
processing any footage:

\begin{lstlisting}[language=bash,caption={SOTA preflight}]
python -m pipeline.sota_registry check
\end{lstlisting}

\noindent If any weight or pin is stale the preflight reports it; resolve before
continuing.  Model weights live in \code{data/comfyui/models/} ($\approx$\,216\,GB,
bind-mounted read-only); do not re-download them.

\subsection{Ingest and frame QA}
\label{sec:pb-ingest}

Place source video files in \code{data/raw/} and submit the job via the web UI
(\texttt{http://localhost:7860}) or directly:

\begin{lstlisting}[language=bash,caption={Step 1 -- ingest and frame extraction}]
python3 -c "
from pipeline.stages import PipelineStages
p = PipelineStages('/data/output/JOB_ID')
result = p.ingest('/data/raw/scene.mp4', fps=4.0)
print(result)
"
\end{lstlisting}

\noindent Ingest extracts frames at 4\,fps to oversample, then the frame-QA
stage selects the sharpest subset.  The QA pipeline applies two complementary
metrics:

\begin{enumerate}
  \item \textbf{MUSIQ no-reference IQA} (\code{pyiqa}, GPU).  A
        perceptual score on the KonIQ-10k scale ($\approx$0--100; good
        photography scores 50--75; severely blurred or noisy footage scores
        below $\approx$25).  Videos unable to supply enough frames above the
        KEEP threshold are marked \code{too\_low\_quality} and skipped.
  \item \textbf{Full-resolution variance-of-Laplacian.}  Sharpness is measured
        at full 4K resolution, not at a downscaled proxy --- downscaling blurs
        away the motion-blur signal (960p blurdetect showed only a
        $\approx$1.6$\times$ range on dreamlab frames; full-4K Laplacian showed a
        $\approx$105$\times$ range).  The \emph{sharpest frame per overlapping
        FIFO window} is kept, which raises the sharpness floor by 20--38\,\% at
        the bottom quartile --- the frames that most poison COLMAP matches.
\end{enumerate}

\noindent Inspect the selected frame count before proceeding:

\begin{lstlisting}[language=bash,caption={Step 3 -- select best frames}]
python3 -c "
from pipeline.stages import PipelineStages
p = PipelineStages('/data/output/JOB_ID')
result = p.select_frames('/data/output/JOB_ID/frames/', target=80)
print(result)
"
# Target: 60-80 frames. If < 40 are above threshold, consider recapture.
\end{lstlisting}

\subsection{Structure-from-Motion and 3DGS training}
\label{sec:pb-sfm-train}

\begin{lstlisting}[language=bash,caption={Step 4 -- COLMAP SfM (sequential matcher for video)}]
python3 -c "
from pipeline.stages import PipelineStages
p = PipelineStages('/data/output/JOB_ID')
result = p.reconstruct('/data/output/JOB_ID/frames_selected/',
                       matcher='sequential')
print(result)
"
# Quality gate: >= 70 % of frames must register.
# If < 50 % register, re-run with matcher='exhaustive'.
\end{lstlisting}

\begin{lstlisting}[language=bash,caption={Step 5 -- 3DGS training (30\,000 iterations)}]
python3 -c "
from pipeline.stages import PipelineStages
p = PipelineStages('/data/output/JOB_ID')
result = p.train('/data/output/JOB_ID/colmap/undistorted/',
                 iterations=30000)
print(result)
"
# Quality gate: PLY file should be > 10 MB; loss < 0.02.
\end{lstlisting}

\noindent The output splat is \code{output/JOB\_ID/model/splat\_30000.ply}.  This
file is the primary LichtFeld-native deliverable that feeds both the Gaussian
splat room path and the object segmentation branch.

\subsection{The stage sequence}
\label{sec:pb-sequence}

The full pipeline in order:

\begin{center}
\begin{tabular}{@{}lL{4.5cm}L{6cm}@{}}
\toprule
\textbf{Stage} & \textbf{What runs} & \textbf{Quality check} \\
\midrule
1 & Ingest: frame extraction at 4\,fps & \code{frames/} directory populated \\
2 & Remove people (if present) & Inspect sample frames \\
3 & Frame-QA selection & 60--80 KEEP-grade frames \\
4 & COLMAP SfM (sequential matcher) & $\geq$\,70\,\% registration \\
5 & 3DGS training (30\,k iterations) & PLY $>$10\,MB, loss $<$0.02 \\
6 & SAM3 concept segmentation & $\geq$\,2 objects identified \\
7 & Per-object PLY extraction + orbit renders & $\geq$5\,k vertices per object \\
7c & TRELLIS.2 hull\_e2e (textured GLB) & Textured GLB produced per object \\
8 & Room mesh extraction (TSDF or 2DGS) & $\geq$\,30\,k vertices, clean topology \\
8b & Mesh cleanup (\code{mesh\_cleaner smooth=0}) & No stray bridging triangles \\
8c & Texture bake (xatlas UV + vertex colours) & Texture atlas per mesh \\
9 & FBX export (\code{blender\_obj\_to\_fbx}) & FBX files produced \\
9b & UE 5.8 import (Nanite game assets) & Assets visible in Content Browser \\
10 & Validate + job complete & Full scene: room + objects \\
\bottomrule
\end{tabular}
\end{center}

\noindent \textbf{Do not stop after training.}  The pipeline is only complete
when it reaches the textured UE scene; stopping at the splat leaves the
deliverable unfinished.

% ─────────────────────────────────────────────────────────────────────
\section{Per-Capture Recipe Selection}
\label{sec:pb-recipe}
% ─────────────────────────────────────────────────────────────────────

Vitrine is not a fixed pipeline.  Each capture is \emph{diagnosed} first; the
router selects the configuration that best fits the capture's bottleneck.  The
four diagnostic axes are:

\begin{description}
  \item[Sharpness] Measured by MUSIQ + full-resolution Laplacian.  Is the
    footage blur-limited?
  \item[Coverage] Measured by a per-surface-point viewpoint count.  Are there
    under-observed regions (holes) or is coverage uniformly dense?
  \item[Hardware] RGB camera only, depth sensor (iPhone LiDAR), or event camera?
  \item[Deliverable] Fidelity-first splat room, polygonal mesh room, or both?
\end{description}

\noindent Table~\ref{tab:pb-recipes} maps common capture profiles to their
recommended configurations.

\begin{table}[ht]
\centering
\caption{Capture profile to recipe selector.  Objects are always reconstructed
via TRELLIS.2 \code{hull\_e2e} regardless of room profile.}
\label{tab:pb-recipes}
\begin{tabular}{@{}L{3.4cm}L{2.4cm}L{2.4cm}L{2.2cm}L{3.0cm}@{}}
\toprule
\textbf{Capture profile} & \textbf{Room rep} & \textbf{Enhancer} & \textbf{Mesh backend} & \textbf{Note} \\
\midrule
\textbf{Blurry $\cdot$ dense $\cdot$ RGB}
  (e.g.\ dreamlab)
  & Splat hero + mesh asset (dual)
  & \emph{None} (ArtiFixer won't help)
  & 2DGS / PGSR
  & Recapture is the real ceiling \\
\addlinespace
\textbf{Sharp $\cdot$ holey $\cdot$ RGB}
  & Splat (primary)
  & \textbf{ArtiFixer} (fills floaters/holes)
  & 2DGS / PGSR if mesh required
  & ArtiFixer's sweet spot \\
\addlinespace
\textbf{Sharp $\cdot$ dense $\cdot$ RGB}
  & Splat (primary)
  & None
  & 2DGS if mesh needed
  & Best case; least pain \\
\addlinespace
\textbf{iPhone LiDAR depth}
  & Mesh (primary)
  & Optional
  & \textbf{AGS-Mesh} (sensor depth)
  & Depth sensor unlocks clean walls \\
\addlinespace
\textbf{Event camera $\cdot$ blurry}
  & Splat
  & ---
  & ---
  & EBAD-Gaussian deblurs at source \\
\bottomrule
\end{tabular}
\end{table}

\subsection{Splat vs.\ mesh for the room}
\label{sec:pb-splat-vs-mesh}

\textbf{Default --- room as splat (NanoGS in UE 5.8).}  For an indoor room the
recommended default is to represent the room as a real Gaussian splat rendered
by the NanoGS UE 5.8 plugin, with object FBXs from TRELLIS.2 embedded
alongside it.  This is simultaneously the highest-fidelity path and the one
that leans most directly on the vendored LichtFeld trainer: the native
\code{splat\_30000.ply} is cleaned with SuperSplat, imported into UE 5.8, and
rendered as real Gaussians.  The mandated polygonal mesh is satisfied by the
TRELLIS.2 object FBXs.

\textbf{When to use a polygonal room mesh.}  If the brief requires a fully
polygonal game-asset scene (Nanite room mesh, not a splat), use \textbf{2DGS}
(quick win, same Open3D-TSDF backend as the current fallback, but 2D-surfel
depth is much cleaner than ellipsoid depth on flat walls) or \textbf{PGSR}
(planar Gaussians, explicitly handles textureless walls, strongest for indoor
rooms).  The legacy CoMe/Open3D-TSDF combination is retained as a fallback
but produces visibly blockier results and should not be the default.

\textbf{When ArtiFixer helps.}  ArtiFixer is a 3DGS refiner targeting
floaters, holes, and under-observed regions: it applies an opacity-mix
enhancement that fills in areas seen from too few viewpoints.  It runs on a
single 48\,GB RTX 6000 Ada (peak $\approx$\,45.5\,GB) and is the correct tool
for \emph{sharp but coverage-limited} captures.  It does \emph{not} deblur.
Applying ArtiFixer to a blur-limited capture wastes GPU time without improving
results.

\textbf{NanoGS floaters.}  The NanoGS plugin renders the splat faithfully,
including any floater Gaussians in the source PLY.  Pre-clean the PLY with
SuperSplat (floater removal by opacity and scale thresholds) before import.
If floaters persist after SuperSplat, check whether coverage was the issue
(ArtiFixer) or whether the source footage was simply blurry (recapture).

\subsection{Object reconstruction}
\label{sec:pb-objects}

Objects follow the same path regardless of room profile:

\begin{enumerate}
  \item \textbf{SAM3 concept segmentation} identifies freestanding objects in
        the splat and extracts per-object Gaussian PLY subsets.
  \item \textbf{Orbit render}: four synthetic camera views around the object
        centroid are rendered from the Gaussian subset.
  \item \textbf{TRELLIS.2 \code{hull\_e2e}}: the orbit renders are submitted to
        TRELLIS.2 (via ComfyUI at \code{http://vitrine-comfyui:8188}) which runs
        FLUX.2 turnaround generation followed by TRELLIS.2 3D reconstruction,
        producing a textured GLB with PBR materials.
  \item \textbf{Prescale}: Blender prescales the GLB to physical units to match
        the COLMAP coordinate frame.
  \item \textbf{FBX export}: the prescaled mesh is exported as FBX for UE 5.8
        game-asset import (Nanite-ready).
\end{enumerate}

\noindent TRELLIS.2 is the validated primary path; Hunyuan3D-2.1 (also staged
in ComfyUI) is the fallback for cases where the orbit renders are cluttered or
partially occluded.  Restart ComfyUI between objects if VRAM is under pressure:
FLUX.2 and TRELLIS.2 cannot co-reside with other large models on 48\,GB.

\subsection{VRAM phasing}
\label{sec:pb-vram}

The pipeline phases large models to avoid out-of-memory conditions.  The key
constraint is that FLUX.2 and TRELLIS.2 together require $\approx$\,45\,GB; they
cannot share VRAM with other resident models.  The orchestrator manages the
lifecycle automatically (load $\to$ infer $\to$ free) but the operator should
be aware:

\begin{itemize}
  \item If ComfyUI reports an OOM, the previous model was not freed cleanly.
        Restart the \code{vitrine-comfyui} container and re-run from the failing
        object.
  \item Process objects serially, not in parallel, unless a second GPU is
        dedicated to the task.
  \item ArtiFixer peaks at $\approx$\,45.5\,GB; do not run it concurrently with
        any ComfyUI workflow.
\end{itemize}

% ─────────────────────────────────────────────────────────────────────
\section{Troubleshooting}
\label{sec:pb-troubleshoot}
% ─────────────────────────────────────────────────────────────────────

\subsection{Recapture flag: footage too soft}
\label{sec:pb-ts-recapture}

\begin{technote}
\textbf{Symptom.}  The frame-QA stage returns \code{VideoQualityVerdict.too\_low\_quality} for all
or most input videos.  The MUSIQ median is below $\approx$\,25.
\end{technote}

The QA gate cannot manufacture sharpness that was never recorded.  Lowering
the MUSIQ threshold feeds more frames to COLMAP but does not improve the
reconstruction --- it just postpones a failure that COLMAP will surface as low
registration or that 3DGS training will surface as a blurry, noisy splat.

\textbf{Resolution: recapture} following the camera settings and movement
guidance in \S\ref{sec:pb-camera-settings}--\ref{sec:pb-movement}.  Checklist:
\begin{itemize}[nosep]
  \item Switch to 60\,fps (halves exposure time vs.\ 30\,fps).
  \item Lock shutter at $\geq$\,1/500\,s in a pro camera app.
  \item Reduce movement speed; add micro-pauses at key viewpoints.
  \item Use a gimbal or brace the camera firmly.
  \item Ensure bright, even illumination (adequate light allows fast shutter
        without under-exposure).
\end{itemize}

\subsection{Directional blur gate}
\label{sec:pb-ts-blur-gate}

Vitrine applies a directional blur metric in addition to the MUSIQ score:
optical-flow magnitude and direction identify frames where camera motion was
predominantly linear (walking forward) versus rotational (panning).  Panning
motion at low shutter speeds produces pronounced horizontal or vertical blur
streaks that fool the Laplacian sharpness metric (which measures spatial
frequency without direction).  Frames with high directional blur are downgraded
before FIFO windowed selection.

If the frame-QA gate discards a high fraction of frames despite the footage
appearing ``mostly sharp'', check whether long pan segments are responsible:
inspect the flow field overlays saved to \code{output/JOB\_ID/qa\_debug/}.
Recapture with slower, more deliberate panning ($<$\,15°/s) or replace panning
segments with lateral translation.

\subsection{COLMAP: low registration rate}
\label{sec:pb-ts-colmap}

A COLMAP registration below 70\,\% usually indicates one of the following:

\begin{description}
  \item[Insufficient overlap] You moved too fast or the scene was only captured
    from one direction.  Recapture with slower movement and the systematic
    multi-pass patterns described in \S\ref{sec:pb-patterns}.
  \item[COLMAP invocation flags] With COLMAP 4.1.0 the option namespace changed;
    old recipes silently use wrong flags.  Use \code{--FeatureExtraction.*} and
    \code{--FeatureMatching.*} (not the old \code{--SiftExtraction.*} namespace).
    Run feature extraction on GPU (\code{--FeatureExtraction.type SIFT} on GPU):
    GPU SIFT achieved $\approx$\,100\,\% registration on the multi-sweep
    dreamlab captures.
  \item[Missing output directory] COLMAP silently fails to create its database if
    the output directory does not already exist.  Run
    \code{mkdir -p colmap/} before invoking \code{feature\_extractor}.
  \item[Mixed cameras] If the capture combined multiple devices, pass
    \code{--ImageReader.single\_camera\_per\_folder} so COLMAP does not force
    one shared intrinsic model across devices.
  \item[Featureless surfaces] Large white walls or uniform floors provide no
    feature points.  Add temporary texture to blank surfaces before
    recapturing (\S\ref{sec:pb-scene-prep}).
\end{description}

\subsection{NanoGS floaters in the UE scene}
\label{sec:pb-ts-floaters}

\begin{description}
  \item[Cause: coverage holes.]  Regions seen from fewer than five viewpoints
    generate noisy or un-anchored Gaussians.  Apply ArtiFixer (the correct tool
    for this profile) before running SuperSplat.
  \item[Cause: motion blur in source footage.]  A blurry Gaussian is still a
    Gaussian; NanoGS renders it faithfully as a blurry blob.  ArtiFixer does
    not help here.  The resolution is recapture.
  \item[Prune with SuperSplat before UE import.] SuperSplat's opacity and scale
    thresholds reliably remove stray low-opacity Gaussians that accumulate in
    poorly observed regions.  Run SuperSplat on every PLY before importing into
    UE regardless of whether ArtiFixer was applied.
\end{description}

\subsection{UE import scale and object placement}
\label{sec:pb-ts-ue-scale}

UE 5.8 imports FBX in centimetres by default; the COLMAP/3DGS coordinate
frame is in metres.  Vitrine's \code{blender\_prescale\_objects} step converts
object meshes to physical scale before FBX export.  If imported objects appear
100$\times$ too large or too small, check that \code{blender\_prescale\_objects}
ran successfully and that the UE import dialogue has \textit{Import Uniform
Scale} set to 1.0 (not 0.01 or 100).  Use \code{ue\_place\_prescaled} (the
keeper placement script) rather than hand-placing objects via Remote Control,
which has been observed to produce incorrect scale on the final assembly step.

\subsection{GPU VRAM OOM during ComfyUI object processing}
\label{sec:pb-ts-vram}

FLUX.2 and TRELLIS.2 together peak at $\approx$\,45.5\,GB.  If a workflow
fails with an out-of-memory error:
\begin{enumerate}[nosep]
  \item Restart the \code{vitrine-comfyui} container to flush all resident
        models: \code{docker restart vitrine-comfyui}.
  \item Re-submit the failing object workflow.  ComfyUI's queue is persistent
        across restarts; the job will resume from the failed node.
  \item If OOM recurs, check that ArtiFixer is not concurrently resident.
        ArtiFixer should be run as a separate stage with ComfyUI idle.
  \item If the \code{vitrine-comfyui} container is sharing GPU\,0 with
        \code{gaussian-toolkit} 3DGS training, defer object processing until
        training is complete.
\end{enumerate}

\noindent The \code{--local\_attn\_size} flag in ArtiFixer is the dominant
VRAM lever: reducing it from the default to a smaller window size (e.g.\ 128)
significantly reduces peak usage at a modest quality cost.  See
\code{docker/artifixer/patches/} for the confirmed fix applied to the
dreamlab run.

% ─────────────────────────────────────────────────────────────────────
\section{Pre-Capture Checklist}
\label{sec:pb-checklist}
% ─────────────────────────────────────────────────────────────────────

Print and check before every capture session.

\begin{center}
\fbox{\begin{minipage}{0.92\textwidth}
\textbf{Pre-capture}\\[2pt]
$\square$\ Scene is fully static: no people, pets, fans, active screens, or flowing water\\
$\square$\ Lighting is even and consistent (meter swing $<$\,1.5 stops over 360°)\\
$\square$\ Mirrors and reflective surfaces covered or removed\\
$\square$\ Featureless walls have temporary texture added\\
$\square$\ Camera battery charged; storage space verified\\[6pt]
\textbf{Camera settings}\\[2pt]
$\square$\ Resolution: 4K; frame rate: 60\,fps\\
$\square$\ Focus: MANUAL, mid-distance\\
$\square$\ Exposure: MANUAL, locked for average scene brightness\\
$\square$\ Shutter: $\geq$\,1/500\,s (locked in pro app)\\
$\square$\ White balance: MANUAL preset matched to light source\\
$\square$\ Lens: 24--50\,mm equivalent (1$\times$ main camera on phones)\\[6pt]
\textbf{Capture}\\[2pt]
$\square$\ Using gimbal or braced two-handed stabilisation\\
$\square$\ Movement speed: $\leq$\,0.5\,m/s; micro-pauses at key viewpoints\\
$\square$\ Perimeter passes at three heights (low / eye / high)\\
$\square$\ Centre diagonal cross-passes\\
$\square$\ Dedicated 360° orbit of each key object at three heights\\
$\square$\ Detail pass of important features (30--50\,cm, 10--15\,s each)\\
$\square$\ Total duration: 60--120\,s per room; 30--60\,s per object\\[6pt]
\textbf{Post-capture}\\[2pt]
$\square$\ Transfer original file via USB or AirDrop (not cloud)\\
$\square$\ No re-encoding unless unavoidable; if required: H.264 CRF\,18, original resolution\\
$\square$\ File placed in \code{data/raw/} on the pipeline host\\
\end{minipage}}
\end{center}
